\NormalFont\ShortTitle{Mateus}
{\MT Evangelho Segundo Mateus

\par }\ChapOne{1}{\PP \VerseOne{1}Livro da geração de Jesus Cristo, filho de Davi, filho de Abraão.
\VS{2}Abraão gerou a Isaque; e Isaque gerou a Jacó; e Jacó gerou a Judá e a seus irmãos.
\VS{3}E Judá gerou de Tamar a Perez e a Zerá; e Perez gerou a Esrom, e Esrom gerou a Arão.
\VS{4}E Arão gerou a Aminadabe; e Aminadabe gerou a Naassom; e Naassom gerou a Salmom.
\VS{5}E Salmom gerou de Raabe a Boaz; e Boaz gerou de Rute a Obede; e Obede gerou a Jessé.
\VS{6}E Jessé gerou ao rei Davi; e o rei Davi gerou Salomão da que
{\ADD{fora mulher}} de Urias.
\VS{7}E Salomão gerou a Roboão; e Roboão gerou a Abias; e Abias gerou a Asa.
\VS{8}E Asa gerou a Josafá; e Josafá gerou a Jorão; e Jorão gerou a Uzias.
\VS{9}E Uzias gerou a Jotão; e Jotão gerou a Acaz; e Acaz gerou a Ezequias.
\VS{10}E Ezequias gerou a Manassés; e Manassés gerou a Amom; e Amom gerou a Josias.
\VS{11}E Josias gerou a Jeconias, e a seus irmãos no
{\ADD{tempo do}} exílio babilônico.
\VS{12}E depois do exílio babilônico Jeconias gerou a Salatiel; e Salatiel gerou a Zorobabel.
\VS{13}E Zorobabel gerou a Abiúde; e Abiúde gerou a Eliaquim; e Eliaquim gerou a Azor.
\VS{14}E Azor gerou a Sadoque; e Sadoque gerou a Aquim; e Aquim gerou a Eliúde.
\VS{15}E Eliúde gerou a Eleazar; e Eleazar gerou a Matã; e Matã gerou a Jacó.
\VS{16}E Jacó gerou a José, o marido de Maria, da qual nasceu Jesus, chamado o Cristo.
\VS{17}De maneira que todas as gerações desde Abraão até Davi são catorze gerações; e desde Davi até o exílio babilônico catorze gerações; e desde o exílio babilônico até Cristo catorze gerações.
\VS{18}E o nascimento de Jesus Cristo foi assim: estando sua mãe Maria desposada com José, antes que se ajuntassem, ela foi achada grávida do Espírito Santo.
\VS{19}Então José, seu marido, sendo justo, e não querendo a expor à infâmia, pensou em deixá-la secretamente.
\VS{20}E ele, pretendendo isto, eis que um anjo do Senhor lhe apareceu em sonho, dizendo: José, filho de Davi, não temas receber Maria, tua mulher, porque o que nela está concebido é do Espírito Santo.
\VS{21}E ela dará à luz um filho, e tu chamarás seu nome Jesus; porque ele salvará seu povo de seus pecados.
\FTNT{*}{{\FR 1:21 }
Jesus significa “O SENHOR salva”
}
\VS{22}E tudo isto aconteceu para que se cumprisse o que foi dito pelo Senhor por meio do profeta, que disse:
\VS{23}Eis que a virgem conceberá, e dará à luz um filho, e chamarão seu nome Emanuel, que traduzido é: Deus conosco.
{\RQ{Isaías 7:14
}}
\VS{24}E despertando José do sonho, fez como o anjo do Senhor tinha lhe mandado, e recebeu sua mulher.
\VS{25}E ele não a conheceu
{\ADD{intimamente}} ,até que ela deu à luz o filho dela, o primogênito, e lhe pôs por nome JESUS.

\par }\Chap{2}{\PP \VerseOne{1}E sendo Jesus já nascido em Belém da Judeia, nos dias do rei Herodes, eis que vieram uns magos do oriente a Jerusalém,
\VS{2}Dizendo: Onde está o Rei nascido dos Judeus? Porque vimos sua estrela no oriente,
\FTNT{*}{{\FR 2:2 }
no oriente
trad.alt. quando ela surgiu – também no v. 9
} e viemos adorá-lo.
\VS{3}E o rei Herodes, ao ouvir
{\ADD{isto}} , ficou perturbado, e com ele toda Jerusalém.
\VS{4}E tendo reunido todos os chefes dos sacerdotes e escribas do povo, perguntou-lhes onde o Cristo havia de nascer.
\VS{5}E eles lhe disseram: Em Belém da Judeia, porque assim está escrito pelo profeta:
\VS{6}E tu Belém, terra de Judá, de maneira nenhuma és a menor entre as lideranças de Judá, porque de ti sairá o Guia que apascentará meu povo Israel.
{\RQ{Miqueias 5:2
}}
\VS{7}Então Herodes, chamando secretamente os magos, perguntou-lhes com precisão sobre o tempo em que a estrela havia aparecido.
\VS{8}E enviando-os a Belém, disse: Ide, e investigai cuidadosamente pelo menino; e quando o achardes, avisai-me, para que também eu venha e o adore.
\VS{9}Depois de ouvirem o rei, eles foram embora. E eis que a estrela que tinham visto no oriente ia adiante deles, até que ela chegou, e ficou parada sobre onde o menino estava.
\VS{10}E eles, vendo a estrela, jubilaram muito com grande alegria.
\VS{11}E entrando na casa, acharam o menino com sua mãe Maria, e prostrando-se o adoraram. E abrindo seus tesouros, ofereceram-lhe presentes: ouro, incenso, e mirra.
\VS{12}E sendo por divina revelação avisados em sonho que não voltassem a Herodes, partiram para sua terra por outro caminho.
\VS{13}E tendo eles partido, eis que um anjo do Senhor apareceu a José em sonho, dizendo: Levanta-te, toma o menino e sua mãe, e foge para o Egito; e fica lá até que eu te diga, porque Herodes buscará o menino para o matar.
\VS{14}Então ele se despertou, tomou o menino e sua mãe de noite, e foi para o Egito;
\VS{15}E esteve lá até a morte de Herodes, para que se cumprisse o que foi dito pelo Senhor por meio do profeta, que disse: Do Egito chamei o meu Filho.
{\RQ{Oseias 11:1
}}
\VS{16}Então Herodes, ao ver que tinha sido enganado pelos magos, irou-se muito, e mandou matar todos os meninos em Belém e em todos os limites de sua região,
{\ADD{da idade}} de dois anos e abaixo, conforme o tempo que tinha perguntado com precisão dos magos.
\VS{17}Então se cumpriu o que foi dito pelo profeta Jeremias, que disse:
\VS{18}Uma voz se ouviu em Ramá, lamentação, choro, e grande pranto; Raquel chorava por seus filhos, e não quis ser consolada, pois já não existem.
{\RQ{Jeremias 31:15
}}
\VS{19}Mas depois de Herodes ter morrido, eis que um anjo do Senhor apareceu no Egito a José em sonho,
\VS{20}Dizendo: Levanta-te, toma o menino e sua mãe, e vai para a terra de Israel, porque já morreram os que procuravam a morte do menino.
\VS{21}Então ele se levantou, tomou o menino e sua mãe, e veio para a terra de Israel.
\VS{22}Porém ao ouvir que Arquelau reinava na Judeia em lugar de seu pai Herodes, ele teve medo de ir para lá; mas avisado por divina revelação em sonho, foi para a região da Galileia,
\VS{23}E veio a habitar na cidade chamada Nazaré, para que se cumprisse o que foi dito pelos profetas, que: Ele será chamado de Nazareno.

\par }\Chap{3}{\PP \VerseOne{1}E naqueles dias veio João Batista, pregando no deserto da Judeia,
\VS{2}E dizendo: Arrependei-vos, porque perto está o Reino dos céus.
\VS{3}Porque este é aquele que foi declarado pelo profeta Isaías, que disse: Voz do que clama no deserto: “Preparai o caminho do Senhor; endireitai suas veredas”.
{\RQ{Isaías 40:3
}}
\VS{4}Este João tinha sua roupa de pelos de camelo e um cinto de couro ao redor de sua cintura, e seu alimento era gafanhotos e mel silvestre.
\VS{5}Então vinham até ele
{\ADD{moradores}} de Jerusalém, de toda a Judeia, e de toda a região próxima do Jordão;
\VS{6}E eram por ele batizados no Jordão, confessando os seus pecados.
\VS{7}Mas quando ele viu muitos dos fariseus e dos saduceus que vinham ao seu batismo, disse-lhes: Ninhada de víboras! Quem vos ensinou a fugir da ira futura?
\VS{8}Dai, pois, frutos condizentes com o arrependimento.
\FTNT{*}{{\FR 3:8 }
arrependimento não é somente um pesar ou remorso pelo pecado, mas sim uma mudança de atitude, abandonando o que é errado
}
\VS{9}E não imagineis, dizendo em vós mesmos: “Temos por pai a Abraão”, porque eu vos digo que até destas pedras Deus pode fazer surgir filhos a Abraão.
\VS{10}E agora mesmoo machado está posto à raiz das árvores; portanto toda árvore que não dá bom fruto é cortada e lançada ao fogo.
\VS{11}Realmente eu vos batizo com água para arrependimento, mas aquele que vem após mim é mais poderoso do que eu; suas sandálias não sou digno de levar. Ele vos batizará com o Espírito Santo e com fogo.
\VS{12}Ele tem sua pá na mão; limpará sua eira, e recolherá seu trigo no celeiro; mas queimará a palha com fogo que nunca se apaga.
\VS{13}Então Jesus veio da Galileia ao Jordão até João para ser por ele batizado.
\VS{14}Mas João lhe impedia, dizendo: Eu preciso ser batizado por ti, e tu vens a mim?
\VS{15}Porém Jesus lhe respondeu: Permite por agora, porque assim nos convém cumprir toda a justiça. Então ele o permitiu.
\VS{16}E tendo Jesus sido batizado, subiu logo da água. E eis que os céus se lhe abriram, e ele viu o Espírito de Deus descendo como uma pomba, vindo sobre ele.
\VS{17}E eis uma voz dos céus, dizendo: Este é o meu Filho amado, em quem me agrado.

\par }\Chap{4}{\PP \VerseOne{1}Então Jesus foi levado pelo Espírito ao deserto para ser tentado pelo diabo.
\VS{2}E depois de jejuar quarenta dias e quarenta noites, teve fome.
\VS{3}E o tentador se aproximou dele, e disse: Se tu és o Filho de Deus, dize que estas pedras se tornem pães.
\VS{4}Mas
{\ADD{Jesus}} respondeu: Está escrito: Não só de pão viverá o ser humano, mas de toda palavra que sai da boca de Deus.
{\RQ{Deuteronômio 8:3
}}
\VS{5}Então o diabo o levou consigo à santa cidade, e o pôs sobre o ponto mais alto do Templo,
\VS{6}E disse-lhe: Se tu és o Filho de Deus, lança-te abaixo, porque está escrito que: Mandará a seus anjos acerca de ti, e te tomarão pelas mãos, para que nunca com teu pé tropeces em pedra alguma.
{\RQ{Salmos 91:11,12
}}
\VS{7}Jesus lhe disse: Também está escrito: Não tentarás
\FTNT{*}{{\FR 4:7 }
tentarás
ou: Não testarás o Senhor teu Deus
} o Senhor teu Deus.
{\RQ{Deuteronômio 6:16
}}
\VS{8}Outra vez o diabo o levou consigo a um monte muito alto, e lhe mostrou todos os reinos do mundo, e a glória deles,
\VS{9}E disse-lhe: Tudo isto te darei se, prostrado, me adorares.
\VS{10}Então Jesus disse: Vai embora, Satanás! Porque está escrito: Ao Senhor teu Deus adorarás, e só a ele cultuarás.
{\RQ{Deuteronômio 6:13
}}
\VS{11}Então o diabo o deixou; e eis que chegaram anjos, e o serviram.
\VS{12}Mas quando Jesus ouviu que João estava preso, voltou para a Galileia.
\VS{13}E deixando Nazaré, veio a morar em Cafarnaum,
{\ADD{cidade}} marítima, nos limites de Zebulom e Naftali,
\VS{14}para que se cumprisse o que foi anunciado pelo profeta Isaías, que disse:
\VS{15}A terra de Zebulom e a terra de Naftali, caminho do mar, além do Jordão, a Galileia dos gentios;
\VS{16}O povo sentado em trevas viu uma grande luz; aos sentados em região e sombra da morte, a luz lhes apareceu.
{\RQ{Isaías 9:1
}}
\VS{17}Desde então Jesus começou a pregar e a dizer: Arrependei-vos, porque perto está o Reino dos céus.
\VS{18}Enquanto Jesus andava junto ao mar da Galileia, viu dois irmãos: Simão, chamado Pedro, e seu irmão André, lançarem a rede ao mar, porque eram pescadores.
\VS{19}E disse-lhes: Vinde após mim, e eu vos farei pescadores de gente.
\VS{20}Então eles logo deixaram as redes e o seguiram.
\VS{21}E passando dali, viu outros dois irmãos: Tiago,
{\ADD{filho}} de Zebedeu, e seu irmão João, em um barco, com seu pai Zebedeu, que estavam consertando suas redes; e ele os chamou.
\VS{22}E eles logo deixaram o barco e seu pai, e o seguiram.
\VS{23}E Jesus rodeava toda a Galileia, ensinando em suas sinagogas, pregando o Evangelho do Reino, e curando toda enfermidade e toda doença no povo.
\VS{24}Sua fama corria por toda a Síria, e traziam-lhe todos que sofriam de algum mal, tendo diversas enfermidades e tormentos, e os endemoninhados, epiléticos, e paralíticos; e ele os curava.
\VS{25}E muitas multidões da Galileia, de Decápolis, de Jerusalém, da Judeia, e dalém do Jordão o seguiam.

\par }\Chap{5}{\PP \VerseOne{1}E quando
{\ADD{Jesus}} viu as multidões, subiu a um monte; e sentando-se, achegaram-se a ele os seus discípulos.
\VS{2}Então ele abriu sua boca e lhes ensinou, dizendo:
\VS{3}Benditos
\FTNT{*}{{\FR 5:3 }
Neste contexto, os benditos, ou bem-aventurados, são as pessoas que têm ou terão motivos de se alegrarem por causa de uma bênção recebida, ainda que não necessariamente estejam contentes no momento
} são os humildes de espírito,
\FTNT{†}{{\FR 5:3 }
humildes ou pobres de espírito são as pessoas que estão aflitas e reconhecem que dependem de Deus para suas necessidades - Isaías 57:15; 66:2
} porque deles é o Reino dos céus.
\VS{4}Benditos são os que choram, porque eles serão consolados.
\VS{5}Benditos são os mansos, porque eles herdarão a terra.
\VS{6}Benditos são os que têm fome e sede de justiça, porque eles serão saciados.
\VS{7}Benditos são os misericordiosos, porque eles alcançarão misericórdia.
\VS{8}Benditos são os limpos de coração, porque eles verão a Deus.
\VS{9}Benditos são os pacíficos, porque eles serão chamados filhos de Deus.
\VS{10}Benditos são os que sofrem perseguição por causa da justiça, porque deles é o Reino dos céus.
\VS{11}Benditos sois vós, quando vos insultarem, perseguirem, e mentirem, falando contra vós todo mal por minha causa.
\VS{12}Jubilai e alegrai-vos, porque grande é vossa recompensa nos céus; pois assim perseguiram aos profetas que foram antes de vós.
\VS{13}Vós sois o sal da terra; mas se o sal perder seu sabor, com que se salgará? Para nada mais presta, a não ser para se lançar fora, e ser pisado pelas pessoas.
\VS{14}Vós sois a luz do mundo; não se pode esconder uma cidade fundada sobre o monte;
\VS{15}Nem se acende a lâmpada para se pôr debaixo de um cesto, mas sim na luminária, e ilumina a todos quantos estão na casa.
\VS{16}Assim brilhe vossa luz diante das pessoas, para que vejam vossas boas obras, e glorifiquem ao vosso Pai, que está nos céus.
\VS{17}Não penseis que vim para revogar a Lei ou os Profetas; não vim para revogar, mas sim para cumprir.
\VS{18}Porque em verdade vos digo que, até que o céu e a terra passem, nem um jota nem um til passará da Lei até que tudo se cumpra.
\FTNT{‡}{{\FR 5:18 }
jota é uma adaptação da letra grega “iota”
}
\VS{19}Portanto qualquer um que desobedecer a um destes menores mandamentos, e assim ensinar às pessoas, será chamado o menor no Reino dos céus; porém qualquer que
{\ADD{os}} cumprir e ensinar, esse será chamado grande no Reino dos céus.
\VS{20}Porque eu vos digo que se a vossa justiça não for maior que a dos escribas e dos fariseus, de maneira nenhuma entrareis no Reino dos céus.
\VS{21}Ouvistes o que foi dito aos antigos: “Não cometerás homicídio”;
\FTNT{§}{{\FR 5:21 }
homicídio
Ref. Êxodo 20:13
} “mas qualquer um que cometer homicídio será réu do julgamento”.
\VS{22}Porém eu vos digo que qualquer um que se irar contra seu irmão sem razão será réu do julgamento. E qualquer um que disser a seu irmão: “Idiota!” será réu do tribunal.
\FTNT{*}{{\FR 5:22 }
tribunal
Ou “sinédrio”, o tribunal dos israelitas
} E qualquer que
{\ADD{lhe}} disser: “Louco!” será réu do fogo do inferno.
\FTNT{†}{{\FR 5:22 }
idiota
significado provável da palavra aramaica às vezes transliterada “raca”
}
\VS{23}Portanto, se trouxeres tua oferta ao altar, e ali te lembrares que teu irmão tem algo contra ti,
\VS{24}Deixa ali tua oferta diante do altar e vai primeiro reconciliar-te com teu irmão, e então vem oferecer a tua oferta.
\VS{25}Faze acordo depressa com teu adversário, enquanto estás com ele no caminho, para não acontecer que o adversário te entregue ao juiz, e o juiz te entregue ao oficial, e te lancem na prisão.
\VS{26}Em verdade te digo que não sairás dali enquanto não pagares a última moeda.
\FTNT{‡}{{\FR 5:26 }
moeda
Ou: quadrante, uma moeda de pequeno valor
}
\VS{27}Ouvistes o que foi dito aos antigos: “Não adulterarás”.
{\RQ{Êxodo 20:14; Deuteronômio 5:18
}}
\VS{28}Porém eu vos digo que qualquer um que olhar para uma mulher para a cobiçar, em seu coração já adulterou com ela.
\VS{29}Se o teu olho direito te faz pecar,
\FTNT{§}{{\FR 5:29 }
Lit. tropeçar (e cair). Também no v. 30
} arranca-o e lança-o de ti; porque é melhor para ti que um dos teus membros se perca do que o teu corpo todo seja lançado no inferno.
\VS{30}E se a tua mão direita te faz pecar, corta-a e lança-a de ti; porque é melhor para ti que um dos teus membros se perca do que o teu corpo todo seja lançado no inferno.
\VS{31}Também foi dito: “Qualquer um que se divorciar de sua mulher, dê a ela carta de divórcio.
{\RQ{Deuteronômio 24:1
}}
\VS{32}Porém eu vos digo que qualquer um que se divorciar de sua mulher, a menos que seja por causa de pecado sexual, faz com que ela adultere; e qualquer um que se casar com a divorciada comete adultério.
\VS{33}Também ouvistes que foi dito aos antigos: “Não jurarás falsamente”, “mas cumprirás ao Senhor os teus juramentos”.
{\RQ{Levítico 19:12; Números 30:2
}}
\VS{34}Porém eu vos digo que de maneira nenhuma jureis; nem pelo céu, porque é o trono de Deus;
\VS{35}Nem pela terra, porque é o suporte de seus pés; nem por Jerusalém, porque é a cidade do grande Rei.
\VS{36}Nem por tua cabeça jurarás, pois nem sequer um cabelo podes tornar branco ou preto.
\VS{37}Mas seja vosso falar: “sim”, “sim”, “não”, “não”; porque o que disso passa procede do maligno.
\VS{38}Ouvistes o que foi dito: “Olho por olho, e dente por dente”.
{\RQ{Êxodo 21:24; Levítico 24:20; Deuteronômio 19:21
}}
\VS{39}Mas eu vos digo que não resistais a quem for mau; em vez disso, a qualquer um que te bater à tua face direita, apresenta-lhe também a outra.
\VS{40}E ao que quiser disputar contigo, e te tomar tua túnica, deixa-lhe também a capa.
\VS{41}E se qualquer um te obrigar a caminhar uma milha, vai com ele duas.
\VS{42}Dá a quem te pedir; e não te desvies de quem quiser de ti tomar emprestado.
\VS{43}Ouvistes o que foi dito: “Amarás teu próximo”, e “odiarás teu inimigo”.
{\RQ{Levítico 19:18
}}
\VS{44}Porém eu vos digo: amai vossos inimigos, bendizei os que vos maldizem, fazei bem aos que vos odeiam, e orai pelos que vos maltratam e vos perseguem,
\VS{45}Para que sejais filhos de vosso Pai que está nos céus; porque ele faz seu sol sair sobre maus e bons, e chover sobre justos e injustos.
\VS{46}Pois se amardes os que vos amam, que recompensa tereis? Não fazem os publicanos também o mesmo?
\VS{47}E se saudardes somente os vossos irmãos, o que fazeis de mais? Não fazem os publicanos também assim?
\VS{48}Portanto, sede vós perfeitos, assim como vosso Pai que está nos céus é perfeito.

\par }\Chap{6}{\PP \VerseOne{1}Ficai atentos para que não façais vossa esmola diante das pessoas a fim de que sejais vistos por elas; de outra maneira não tereis recompensa de vosso Pai que está nos céus.
\VS{2}Portanto, quando fizeres esmola, não faças tocar trombeta diante de ti, como fazem os hipócritas nas sinagogas e nas ruas, para serem honrados pelas pessoas; em verdade vos digo que já receberam sua recompensa.
\VS{3}Mas quando tu fizeres esmola, não saiba tua mão esquerda o que faz a tua direita;
\VS{4}Para que a tua esmola seja em segredo, e teu Pai, que vê em segredo, ele te recompensará em público.
\VS{5}E quando orares, não sejas como os hipócritas; porque eles amam orar em pé nas sinagogas e nas esquinas das ruas para serem vistos pelas pessoas. Em verdade vos digo que já receberam sua recompensa.
\VS{6}Porém tu, quando orares, entra em teu quarto, fecha tua porta, e ora a teu Pai, que está em segredo; e teu Pai, que vê em segredo, ele te recompensará em público.
\VS{7}E quando orardes, não façais repetições inúteis como os gentios, que pensam que por muito falarem serão ouvidos.
\VS{8}Não sejais, pois, semelhantes a eles; porque vosso Pai sabe o que necessitais, antes que vós peçais a ele.
\VS{9}Vós, portanto, orareis assim: Pai nosso, que estás nos céus, santificado seja o teu nome.
\VS{10}Venha o teu Reino. Seja feita a tua vontade, tanto na terra, assim como no céu.
\VS{11}O pão nosso de cada dia nos dá hoje.
\VS{12}E perdoa-nos nossas dívidas, assim como nós também perdoamos aos nossos devedores.
\VS{13}E não nos conduzas à tentação, mas livra-nos do mal; porque teu é o Reino, o poder, e a glória, para sempre, Amém.
\VS{14}Porque se perdoardes às pessoas suas ofensas, vosso Pai celestial também vos perdoará;
\VS{15}Mas se não perdoardes às pessoas suas ofensas, também vosso Pai não vos perdoará vossas ofensas.
\VS{16}E quando jejuardes, não vos mostreis tristonhos, como os hipócritas; porque eles desfiguram seus rostos, para parecerem aos outros que jejuam. Em verdade vos digo que eles já receberam sua recompensa.
\VS{17}Porém tu, quando jejuares, unge a tua cabeça e lava o teu rosto,
\VS{18}Para não pareceres às pessoas que jejuas, mas sim ao teu Pai, que está em segredo; e teu Pai, que vê em segredo, ele te recompensará em público.
\VS{19}Não ajunteis para vós tesouros na terra, onde e a traça e a ferrugem gastam, e onde os ladrões invadem e roubam;
\VS{20}Mas ajuntai para vós tesouros no céu, onde nem a traça nem a ferrugem gastam, e onde os ladrões não invadem nem roubam.
\VS{21}Porque onde estiver o vosso tesouro, ali estará também o vosso coração.
\VS{22}A lâmpada do corpo é o olho; portanto, se o teu olho for puro,
\FTNT{*}{{\FR 6:22 }
Significando “sincero”, “franco”. A palavra grega significa literalmente “simples”, “sem mistura”, no caso, aplicado literalmente ao olho, seria “sem nada que impeça a passagem da luz”
} todo o teu corpo será cheio de luz.
\VS{23}Porém se o teu olho for maligno, todo o teu corpo será cheio de trevas. Assim, se a luz que há em ti são trevas, como são grandes
{\ADD{essas}} trevas!
\VS{24}Ninguém pode servir a dois senhores; pois ou odiará um e amará outro; ou se apegará a um, e desprezará o outro. Não podeis servir a Deus e às riquezas.
\FTNT{†}{{\FR 6:24 }
riquezas
lit. “Mamom”, a personificação dos bens materiais e da ganância
}
\VS{25}Por isso vos digo: não andeis ansiosos por vossa vida, sobre o que haveis de comer, ou que haveis de beber; nem por vosso corpo, sobre com que vos haveis de vestir. Não é a vida mais que o alimento, e o corpo
{\ADD{mais}} que a roupa?
\VS{26}Olhai para as aves do céu, que não semeiam, nem ceifam, nem ajuntam em celeiros; e contudo vosso Pai celestial as alimenta. Não sois vós muito mais importantes que elas?
\VS{27}E qual de vós poderá, por sua ansiedade, acrescentar um côvado à sua estatura?
\FTNT{‡}{{\FR 6:27 }
estatura
ou: qual de vós poderá, por sua ansiedade, alongar um pouco sua vida?
}
\VS{28}E por que andais ansiosos pela roupa? Prestai atenção aos lírios do campo, como crescem; eles nem trabalham nem fiam.
\VS{29}Mas eu vos digo que nem mesmo Salomão, em toda a sua glória, se vestiu como um deles.
\VS{30}Se Deus veste desta maneira a erva do campo, que hoje existe, e amanhã é lançada no forno, não
{\ADD{vestirá}} ele muito mais a vós, que tendes pouca fé?
\VS{31}Não andeis, pois, ansiosos, dizendo: “Que comeremos?” Ou “Que beberemos?” Ou “Com que nos vestiremos?”
\VS{32}Porque os gentios buscam todas estas coisas, e vosso Pai celestial sabe que necessitais destas coisas, todas elas.
\VS{33}Mas buscai primeiramente o Reino de Deus e a sua justiça, e todas estas coisas vos serão acrescentadas.
\VS{34}Não andeis, pois, ansiosos pelo dia de amanhã, porque o amanhã cuidará de si mesmo. Basta a cada dia o seu mal.

\par }\Chap{7}{\PP \VerseOne{1}Não julgueis, para que não sejais julgados.
\VS{2}Porque com o juízo que julgardes sereis julgados; e com a medida que medirdes vos medirão de volta.
\VS{3}Ora, por que vês o cisco que está no olho de teu irmão e não enxergas a trave que está em teu próprio olho?
\VS{4}Ou como dirás a teu irmão: “Deixa-me tirar o cisco do teu olho”, se eis que há uma trave em teu próprio olho?
\VS{5}Hipócrita! Tira primeiro a trave do teu olho, e então verás claramente para tirar o cisco do olho de teu irmão.
\VS{6}Não deis o que é santo aos cães, nem lanceis vossas pérolas diante dos porcos, para não acontecer de as pisarem com os pés e, virando-se, vos despedacem.
\VS{7}Pedi, e vos será dado; buscai, e achareis; batei, e vos será aberto;
\VS{8}Pois qualquer um que pede recebe; e quem busca acha; e ao que bate lhe é aberto.
\VS{9}E quem há dentre vós que, se seu filho pedir pão, lhe dará uma pedra?
\VS{10}E se pedir peixe, lhe dará uma serpente?
\VS{11}Ora, se vós, sendo maus, sabeis dar bons presentes a vossos filhos, quanto mais o vosso Pai, que está nos céus, dará coisas boas aos que lhe pedirem!
\VS{12}Portanto tudo o que quiserdes que os outros vos façam, fazei-lhes vós também assim; porque esta é a Lei e os Profetas.
\VS{13}Entrai pela porta estreita; porque larga é a porta, e espaçoso o caminho que leva à perdição; e muitos são os que por ela entram.
\VS{14}Porque estreita é a porta, e apertado o caminho que leva à vida; e são poucos os que a acham.
\VS{15}Tende cuidado, porém, com os falsos profetas, que vêm a vós com roupa de ovelhas, mas por dentro são lobos vorazes.
\VS{16}Vós os conhecereis pelos seus frutos. Por acaso se colhem uvas dos espinheiros, ou figos dos cardos?
\VS{17}Assim toda boa árvore dá bons frutos, mas a arvore má dá frutos maus.
\VS{18}A boa árvore não pode dar frutos maus, nem a árvore má dar bons frutos.
\VS{19}Toda árvore que não dá bom fruto é cortada e lançada ao fogo.
\VS{20}Portanto vós os conhecereis pelos seus frutos.
\VS{21}Não é qualquer um que me diz: “Senhor, Senhor” que entrará no Reino dos céus; mas sim aquele que faz a vontade do meu Pai que está nos céus.
\VS{22}Muitos me dirão naquele dia: “Senhor, Senhor! Não profetizamos em teu nome? E em teu nome não expulsamos os demônios? E em teu nome não fizemos muitas maravilhas?”
\VS{23}Então claramente lhes direi: “Nunca vos conheci. Afastai-vos de mim, transgressores!”
\FTNT{*}{{\FR 7:23 }
transgressores
lit. praticantes de transgressão (ou maldade
}
\VS{24}Portanto todo o que ouve estas minhas palavras e as pratica, eu o compararei ao homem prudente, que construiu sua casa sobre a rocha.
\VS{25}E a chuva desceu, correntezas vieram, ventos sopraram, e atingiram aquela casa; e ela não caiu, porque estava fundada sobre a rocha.
\VS{26}Porém todo o que ouve estas minhas palavras e não as pratica, eu o compararei ao homem tolo, que construiu sua casa sobre a areia.
\VS{27}E a chuva desceu, correntezas vieram, ventos sopraram, e atingiram aquela casa; e ela caiu, e sua queda foi grande.
\VS{28}E aconteceu que, quando Jesus terminou estas palavras, as multidões estavam admiradas de sua doutrina,
\VS{29}porque ele os ensinava como tendo autoridade, e não como os escribas.

\par }\Chap{8}{\PP \VerseOne{1}Quando ele desceu do monte, muitas multidões o seguiram.
\VS{2}E eis que um leproso veio e prostrou-se diante dele, dizendo: Senhor, se quiseres, podes me limpar.
\VS{3}Jesus estendeu a mão e o tocou, dizendo: Quero, sê limpo. E logo ele ficou limpo de sua lepra.
\VS{4}Então Jesus lhe disse: Tem o cuidado de dizeres a ninguém; mas vai mostrar-te ao sacerdote, e oferece a oferta que Moisés ordenou, para que lhes haja testemunho.
\VS{5}Quando Jesus entrou em Cafarnaum, veio a ele um centurião, rogando-lhe,
\VS{6}E dizendo: Senhor, o meu servo jaz em casa, paralítico, e gravemente atormentado.
\VS{7}E Jesus lhe disse: Eu irei, e o curarei.
\VS{8}E o centurião respondeu: Senhor, não sou digno de que entres sob meu telhado; mas dize somente uma palavra, e o meu servo sarará.
\VS{9}Pois eu também sou homem debaixo de autoridade, e tenho debaixo de meu comando soldados; e digo a este: “Vai”, e ele vai; e ao outro: “Vem”, e vem; e a meu servo: “Faze isto”, e ele faz.
\VS{10}Quando Jesus ouviu
{\ADD{isto}} , maravilhou-se, e disse aos que o seguiam: Em verdade vos digo que nem mesmo em Israel encontrei tanta fé.
\VS{11}Mas eu vos digo que muitos virão do oriente e do ocidente, e se sentarão à mesa com Abraão, Isaque, e Jacó, no Reino dos céus.
\VS{12}Os filhos do reino, porém, serão lançados nas trevas de fora; ali haverá pranto e ranger de dentes.
\VS{13}Então Jesus disse ao centurião: Vai, e assim como creste, a ti seja feito.E naquela mesma hora o servo dele foi sarado.
\VS{14}E quando Jesus veio à casa de Pedro, viu a sogra dele, deitada e com febre.
\VS{15}Ele tocou a mão dela, e a febre a deixou. Então ela se levantou e começou a servi-los.
\VS{16}Quando chegou o anoitecer, trouxeram-lhe muitos endemoninhados. Ele expulsou-lhes os espíritos com a palavra, e curou todos os que estavam doentes,
\VS{17}Para que se cumprisse o que havia sido dito pelo profeta Isaías, que disse: Ele tomou sobre si as nossas enfermidades, e levou as nossas doenças.
{\RQ{Isaías 53:4
}}
\VS{18}E Jesus, ao ver muitas multidões ao redor de si, mandou que passassem para a outra margem.
\VS{19}Então um escriba se aproximou, e disse-lhe: Mestre, eu te seguirei aonde quer que fores.
\VS{20}Jesus lhe respondeu: As raposas têm covis, e as aves do céu ninhos; mas o Filho do homem não tem onde recostar a cabeça.
\VS{21}E outro dos seus discípulos lhe disse: Senhor, permite-me ir primeiro enterrar meu pai.
\VS{22}Porém Jesus lhe disse: Segue-me, e deixa aos mortos enterrarem seus mortos.
\VS{23}Então ele entrou no barco, e seus discípulos o seguiram.
\VS{24}E eis que se levantou no mar uma tormenta tão grande que o barco era coberto pelas ondas; porém ele dormia.
\VS{25}E seus discípulos se aproximaram para acordá-lo, dizendo: Senhor, salva-nos! Estamos sendo destruídos!
\VS{26}E ele lhes respondeu: Por que temeis,
{\ADD{homens}} de pouca fé? Então ele se levantou e repreendeu os ventos e o mar. E houve grande calmaria.
\VS{27}E aqueles homens se maravilharam, dizendo: Quem é este, que até os ventos e o mar lhe obedecem?
\VS{28}E quando chegou à outra margem, à terra dos gergesenos, vieram-lhe ao encontro dois endemoninhados que tinham saído dos sepulcros. Eles eram tão ferozes que ninguém podia passar por aquele caminho.
\VS{29}E eis que gritaram, dizendo: Que temos contigo, Jesus, Filho de Deus? Vieste aqui nos atormentar antes do tempo?
\VS{30}Enquanto isso longe deles estava uma manada de muitos porcos pastando.
\VS{31}E os demônios rogaram-lhe, dizendo: Se nos expulsares, permite-nos entrar naquela manada de porcos.
\VS{32}E ele lhes disse: Ide. Então eles saíram, e entraram na manada de porcos; e eis que toda aquela manada de porcos se lançou de um precipício ao mar, e morreram nas águas.
\VS{33}E os que cuidavam dos porcos fugiram; e ao chegarem à cidade, anunciaram todas
{\ADD{estas}} coisas, inclusive o que
{\ADD{havia acontecido}} aos endemoninhados.
\VS{34}E eis que toda aquela cidade saiu ao encontro de Jesus; e quando o viram, rogaram-lhe que se retirasse do território deles.

\par }\Chap{9}{\PP \VerseOne{1}Então ele entrou no barco, passou para a outra margem, e veio à sua própria cidade.
\VS{2}E eis que lhe trouxeram um paralítico, deitado em um leito. Quando Jesus viu a fé deles, disse ao paralítico: Tem bom ânimo, filho! Teus pecados são perdoados.
\VS{3}E eis que alguns dos escribas disseram entre si: Ele blasfema.
\VS{4}Mas Jesus, conhecendo seus pensamentos, disse: “Por que pensais o mal em vossos corações?
\VS{5}Pois o que é mais fácil? Dizer: ‘Teus pecados estão perdoados’, ou dizer: ‘Levanta-te, e anda’?
\VS{6}Ora, para que saibais que o Filho do homem tem autoridade na terra para perdoar pecados”, (Ele, então, disse ao paralítico): Levanta-te, toma o teu leito, e vai para tua casa.
\VS{7}E este se levantou e foi para sua casa.
\VS{8}Quando as multidões viram
{\ADD{isto}} , ficaram maravilhadas, e glorificaram a Deus, que tinha dado tal autoridade aos homens.
\VS{9}E Jesus, ao passar dali, viu um homem sentado na coletoria de impostos, chamado Mateus; e disse-lhe: Segue-me. Então este se levantou e o seguiu.
\VS{10}E aconteceu que, enquanto
{\ADD{Jesus}} estava reclinado à mesa na casa
{\ADD{de Mateus}} , eis que muitos publicanos e pecadores vieram e se reclinaram à mesa juntamente com Jesus e seus discípulos.
\FTNT{*}{{\FR 9:10 }
reclinar-se à mesa
i.e., ter uma refeição
}
\VS{11}E quando os fariseus viram
{\ADD{isto}} , perguntaram aos seus discípulos: Por que o vosso Mestre come com publicanos e pecadores?
\VS{12}Porém Jesus ouviu, e respondeu-lhes: Os que têm saúde não precisam de médico, mas sim os que estão doentes.
\VS{13}Mas ide aprender o que significa: ‘Quero misericórdia, e não sacrifício’.
{\RQ{Oseías 6:6
}} Porque eu não vim chamar os justos, mas sim, os pecadores, ao arrependimento.
\VS{14}Então os discípulos de João vieram a ele, e perguntaram: Por que nós e os fariseus jejuamos muitas vezes, mas os teus discípulos não jejuam?
\VS{15}E Jesus lhes respondeu: Podem, por acaso, os convidados do casamento
\FTNT{†}{{\FR 9:15 }
convidados do casamento
lit. filhos da câmara nupcial
} andar tristes enquanto o noivo está com eles? Mas dias virão, quando o noivo lhes for tirado, e então jejuarão.
\VS{16}E ninguém põe remendo de pano novo em roupa velha; porque tal remendo rasga a roupa, e o rompimento se torna pior.
\VS{17}Nem põem vinho novo em odres
\FTNT{‡}{{\FR 9:17 }
odre
recipiente para líquidos de couro. Os odres novos podiam inchar com o resultado da fermentação do vinho novo. Os odres velhos já tinham perdido a capacidade de se esticarem, e por isso não servem para se pôr vinho novo
} velhos; pois senão os odres se rompem, o vinho se derrama, e os odres se perdem; mas põem o vinho novo em odres novos, e ambos juntamente se conservam.
\VS{18}Enquanto ele lhes dizia estas coisas, eis que um chefe
{\ADD{de sinagoga}} veio prostrar-se diante dele, e disse: Minha filha faleceu ainda agora; mas vem, e põe tua mão sobre ela, e ela viverá.
\VS{19}Então Jesus se levantou e o seguiu com seus discípulos.
\VS{20}(Eis, porém, que uma mulher enferma de um fluxo de sangue havia doze anos veio por detrás
{\ADD{dele}} , e tocou a borda de sua roupa;
\VS{21}Porque dizia consigo mesma: Se eu tão-somente tocar a roupa dele, serei curada.
\VS{22}Jesus se virou e a viu. Então disse: Tem bom ânimo, filha, a tua fé te sarou. E desde aquela hora a mulher ficou com saúde.)
\VS{23}Quando Jesus chegou à casa daquele chefe, viu os tocadores de flauta e a multidão que fazia alvoroço,
\VS{24}E disse-lhes: Retirai-vos, porque a menina não está morta, mas sim dormindo. E riram dele.
\VS{25}Mas quando a multidão foi expulsa, ele entrou, pegou a mão dela, e a menina se levantou.
\VS{26}E esta notícia se espalhou por toda aquela terra.
\VS{27}E saindo Jesus dali, dois cegos o seguiram, gritando: Tem compaixão de nós, Filho de Davi!
\VS{28}E quando ele entrou em casa, os cegos vieram a ele. Jesus lhes perguntou: Credes que posso fazer isto? Eles lhe responderam: Sim, Senhor.
\VS{29}Então tocou os olhos deles, dizendo: Seja feito convosco conforme a vossa fé.
\VS{30}E os olhos deles se abriram. Então Jesus os advertiu severamente, dizendo: Tende o cuidado de que ninguém saiba disso.
\VS{31}Porém eles saíram e divulgaram a notícia acerca dele por toda aquela terra.
\VS{32}Enquanto eles saíam, eis que lhe trouxeram um homem mudo e endemoninhado.
\VS{33}Quando o demônio foi expulso, o mudo passou a falar. Então as multidões ficaram maravilhadas, e disseram: Nunca se viu algo assim em Israel!
\VS{34}Mas os fariseus diziam: É pelo chefe dos demônios que ele expulsa os demônios.
\VS{35}Jesus percorria todas as cidades e aldeias, ensinando em suas sinagogas, pregando o Evangelho do Reino, e curando toda enfermidade e toda doença entre o povo.
\VS{36}Quando ele viu as multidões, teve compaixão delas, porque andavam cansadas e desamparadas, como ovelhas que não têm pastor.
\VS{37}Então disse aos seus discípulos: Em verdade a colheita é grande, porém os trabalhadores são poucos.
\VS{38}Portanto rogai ao Senhor da colheita que envie trabalhadores à sua colheita.

\par }\Chap{10}{\PP \VerseOne{1}{\ADD{Jesus}} chamou a si os seus doze discípulos, e deu-lhes poder sobre os espíritos imundos, para os expulsarem, e curarem toda enfermidade e toda doença.
\VS{2}E os nomes dos doze apóstolos são estes: o primeiro, Simão, chamado Pedro, e seu irmão André; Tiago,
{\ADD{filho}} de Zebedeu, e seu irmão João;
\VS{3}Filipe e Bartolomeu; Tomé, e Mateus o publicano; Tiago,
{\ADD{filho}} de Alfeu; e Lebeu, por sobrenomeTadeu;
\VS{4}Simão o zeloso,
\FTNT{*}{{\FR 10:4 }
zeloso
ou “zelote” - tradicionalmente “cananeu”
} e Judas Iscariotes, o mesmo que o traiu.
\FTNT{†}{{\FR 10:4 }
traiu
lit. entregou
}
\VS{5}Jesus enviou esses doze, e lhes mandou, dizendo: Não ireis pelo caminho dos gentios, nem entrareis em cidade de samaritanos.
\VS{6}Em vez disso, ide às ovelhas perdidas da casa de Israel.
\VS{7}E quando fordes, proclamai, dizendo: ‘Perto está o Reino dos céus’.
\VS{8}Curai os doentes, limpai os leprosos, ressuscitai os mortos, expulsai os demônios; recebestes de graça, dai de graça.
\VS{9}Não tomeis convosco ouro, nem prata, nem cobre em vossos cintos;
\VS{10}nem bolsas para o caminho, nem duas túnicas, nem sandálias, nem bordão;
\FTNT{‡}{{\FR 10:10 }
i.e., uma vara ou cajado para facilitar a caminhada
} pois o trabalhador é digno de seu alimento.
\VS{11}E em qualquer cidade ou aldeia que entrardes, informai-vos de quem nela seja digno, e ficai ali até que saiais.
\VS{12}E quando entrardes na casa, saudai-a.
\VS{13}Se a casa for digna, venha sobre ela a vossa paz; mas se ela não for digna, volte para vós a vossa paz.
\VS{14}E quem quer que não vos receber, nem ouvir vossas palavras, quando sairdes daquela casa ou cidade, sacudi o pó de vossos pés.
\VS{15}Em verdade vos digo que no dia do julgamento mais tolerável será para a região de Sodoma e Gomorra do que para aquela cidade.
\VS{16}Eis que eu vos envio como ovelhas em meio de lobos; portanto, sede prudentes como as serpentes e inofensivos como pombas.
\VS{17}Porém tende cuidado com as pessoas; porque vos entregarão em tribunais, e vos açoitarão em suas sinagogas;
\VS{18}E até perante governadores e reis sereis levados por causa de mim, para que haja testemunho a eles e aos gentios.
\VS{19}Mas quando vos entregarem, não estejais ansiosos de como ou que falareis; porque naquela mesma hora vos será dado o que deveis falar.
\VS{20}Porque não sois vós os que falais, mas sim o Espírito do vosso Pai que fala em vós.
\VS{21}E irmão entregará irmão à morte, e pai ao filho; e filhos se levantarão contra os pais, e os matarão.
\VS{22}E sereis odiados por todos por causa de meu nome; mas aquele que perseverar até o fim, esse será salvo.
\VS{23}Quando, então, vos perseguirem nesta cidade, fugi para outra; porque em verdade vos digo que não acabareis de
{\ADD{percorrer}} as cidades de Israel, até que venha o Filho do homem.
\VS{24}O discípulo não é superior ao mestre, nem o servo superior ao seu senhor.
\VS{25}Seja suficiente ao discípulo ser como o seu mestre, e ao servo como o seu senhor; se ao chefe da casa chamaram de Belzebu, quanto mais aos membros de sua casa!
\VS{26}Portanto, não os temais; porque nada há encoberto que não se revelará,
{\ADD{nada}} oculto que não se saberá.
\VS{27}O que eu vos digo em trevas, dizei na luz; e o que ouvis ao ouvido, proclamai sobre os telhados.
\VS{28}E não temais os que matam o corpo, mas não podem matar a alma; temei mais aquele que pode destruir tanto a alma como o corpo no inferno.
\VS{29}Não se vendem dois pardais por uma pequena moeda? Mas nem um deles cairá em terra contra a vontade de
\FTNT{§}{{\FR 10:29 }
contra a vontade de
lit. sem
} vosso Pai.
\VS{30}E até os cabelos de vossas cabeças estão todos contados.
\VS{31}Assim, não tenhais medo; mais valeis vós que muitos pardais.
\VS{32}Portanto, todo aquele que me der reconhecimento diante das pessoas, também eu o reconhecerei diante de meu Pai, que está nos céus.
\VS{33}Porém qualquer um que me negar diante das pessoas, também eu o negarei diante de meu Pai, que está nos céus.
\VS{34}Não penseis que vim trazer paz à terra. Não vim trazer paz, mas sim espada.
\VS{35}Porque eu vim pôr em discórdia ‘o homem contra seu pai, a filha contra sua mãe, e a nora contra sua sogra.
\VS{36}E os inimigos do homem serão os de sua própria casa’.
{\RQ{Miqueias 7:6
}}
\VS{37}Quem ama pai ou mãe mais que a mim não é digno de mim; e quem ama filho ou filha mais que a mim não é digno de mim;
\VS{38}E quem não toma sua cruz e segue após mim não é digno de mim.
\VS{39}Quem achar sua vida a perderá; e quem, por causa de mim, perder sua vida, a achará.
\VS{40}Quem vos recebe, recebe a mim; e quem me recebe, recebe aquele que me enviou.
\VS{41}Quem recebe um profeta por reconhecê-lo como
\FTNT{*}{{\FR 10:41 }
por reconhecê-lo como
lit. em nome de – também v. 42
} profeta receberá recompensa de profeta; e quem recebe um justo por reconhecê-lo como justo receberá recompensa de justo.
\VS{42}E qualquer um que der ainda que somente um copo de
{\ADD{água}} fria a um destes pequenos por reconhecê-lo como discípulo, em verdade vos digo que de maneira nenhuma perderá sua recompensa.

\par }\Chap{11}{\PP \VerseOne{1}Quando Jesus acabou de dar as ordens aos seus doze discípulos, partiu dali para ensinar e para pregar em suas cidades.
\VS{2}E João, ao ouvir na prisão as obras de Cristo, enviou
{\ADD{-lhe}} dois de seus discípulos,
\VS{3}Perguntando-lhe: És tu aquele que havia de vir, ou esperamos outro?
\VS{4}Jesus lhes respondeu: Ide anunciar a João as coisas que ouvis e vedes:
\VS{5}Os cegos veem, e os mancos andam; os leprosos são limpos, e os surdos ouvem; os mortos são ressuscitados, e aos pobres é anunciado o Evangelho;
\VS{6}E bendito
\FTNT{*}{{\FR 11:6 }
bendito
ou “bem-aventurado”
} é aquele que não se ofender
\FTNT{†}{{\FR 11:6 }
não se ofender
Ou: “Perder a fé”. Lit. “não tropeçar”. Tradicionalmente “escandalizar-se”
} em mim.
\VS{7}Depois que eles se foram, Jesus começou a dizer às multidões acerca de João: Que saístes ao deserto para ver? Uma cana que se move pelo vento?
\VS{8}Mas que saístes para ver? Um homem vestido com roupas delicadas? Eis que os que usam roupas delicadas estão nas casas dos reis.
\VS{9}Mas que saístes para ver? Um profeta? Sim, eu vos digo, e muito mais que um profeta;
\VS{10}Porque este é aquele sobre o qual está escrito: ‘Eis que diante de tua face envio o meu mensageiro, que preparará o teu caminho diante de ti.’
{\RQ{Malaquias 3:1
}}
\VS{11}Em verdade vos digo que, dentre os nascidos de mulheres, não se levantou
{\ADD{outro}} maior que João Batista; porém o menor no Reino dos céus é maior que ele.
\VS{12}E desde os dias de João Batista até agora, o Reino dos céus está sujeito à violência, e os que usam de violência se apoderam dele.
\VS{13}Porque todos os profetas e a Lei profetizaram até João.
\VS{14}E se estais dispostos a aceitar, este é o Elias que havia de vir.
\VS{15}Quem tem ouvidos para ouvir, ouça.
\VS{16}Mas com quem compararei esta geração? Semelhante é às crianças que se sentam nas praças, e chamam aos seus colegas,
\VS{17}E dizem: ‘Tocamos flauta para vós, mas não dançastes; cantamos lamentações para vós, mas não chorastes.’
\VS{18}Porque veio João, sem comer nem beber, e dizem: ‘Ele tem demônio.’
\VS{19}Veio o Filho do homem, que come e bebe, e dizem: ‘Eis aqui um homem comilão e beberrão, amigo de publicanos e pecadores!’ Mas a sabedoria prova-se justa por meio de seus filhos.
\VS{20}Então ele começou a repreender as cidades em que a maioria de seus milagres haviam sido feitos, por não terem se arrependido:
\VS{21}Ai de ti Corazim! Ai de ti Betsaida! Porque se em Tiro e em Sidom tivessem sido feitos os milagres que em vós foram feitos, há muito tempo teriam se arrependido com saco e com cinza!
\VS{22}Porém eu vos digo que mais tolerável será para Tiro e Sidom, no dia do juízo, que para vós.
\VS{23}E tu, Cafarnaum, que estás exaltada até o céu, ao Xeol
\FTNT{‡}{{\FR 11:23 }
Xeol é o lugar dos mortos
} serás derrubada! Pois se em Sodoma tivessem sido feitos os milagres que foram feitos em ti, ela teria permanecido até hoje.
\VS{24}Porém eu vos digo que mais tolerável será para os da região de Sodoma, no dia de juízo, que para ti.
\VS{25}Naquele tempo Jesus pronunciou: Graças te dou, Pai, Senhor do céu e da terra, porque escondeste estas coisas aos sábios e entendidos, e as revelaste às crianças.
\VS{26}Sim, Pai, porque assim foi agradável a ti.
\VS{27}Todas as coisas me foram entregues pelo meu Pai; e ninguém conhece o Filho, a não ser o Pai; nem ninguém conhece o Pai, a não ser o Filho, e a quem o Filho
{\ADD{o}} quiser revelar.
\VS{28}Vinde a mim todos vós que estais cansados e sobrecarregados, e eu vos farei descansar.
\VS{29}Tomai sobre vós o meu jugo e aprendei de mim, porque sou manso e humilde de coração; e encontrareis descanso para as vossas almas.
\VS{30}Pois o meu jugo é suave, e minha carga é leve.

\par }\Chap{12}{\PP \VerseOne{1}Naquele tempo Jesus estava indo pelas plantações de cereais no sábado. Seus discípulos tinham fome, e começaram a arrancar espigas e a comer.
\VS{2}Quando os fariseus viram, disseram-lhe: Eis que os teus discípulos fazem o que não é lícito fazer no sábado.
\VS{3}Ele, porém, lhes disse: Não lestes o que Davi fez quando teve fome, ele e os que com ele estavam,
\VS{4}Como ele entrou na casa de Deus, e comeu os pães da proposição, que a ele não era lícito comer, nem também aos que com ele estavam, a não ser somente aos sacerdotes?
\FTNT{*}{{\FR 12:4 }
Os pães da proposição eram oferecidos a Deus semanalmente no templo, e depois eram comidos pelos sacerdotes
Levítico 24:5-9; Êxodo 25:30
}
\VS{5}Ou não lestes na Lei que, nos sábados, os sacerdotes no Templo profanam o sábado, sem se tornarem culpados?
\VS{6}Eu, porém, vos digo que o maior que o Templo está aqui.
\VS{7}Mas se vós soubésseis o que significa: ‘Quero misericórdia, e não sacrifício’,
{\RQ{Oseias 6:6
}} não condenaríeis os inocentes.
\VS{8}Porque o Filho do homem é Senhor até do sábado.
\VS{9}E partindo dali,
{\ADD{Jesus}} entrou na sinagoga deles.
\VS{10}E eis que havia ali um homem que tinha uma mão definhada; e eles, a fim de o acusarem, perguntaram-lhe: É lícito curar nos sábados?
\VS{11}E ele lhes respondeu: Qual de vós será a pessoa que, caso tenha uma ovelha, e se a tal cair em uma cova no sábado, não usará de sua força para a levantar?
\VS{12}Ora, quanto mais vale um ser humano que uma ovelha! Assim, pois, é lícito fazer o bem nos sábados.
\VS{13}Então disse para aquele homem: Estende a tua mão. Ele a estendeu, e foi-lhe restaurada, sã como a outra.
\VS{14}Então os fariseus saíram e se reuniram para planejar contra ele, como o matariam.
\VS{15}Mas Jesus, sabendo disso, retirou-se dali. Muitas multidões o seguiram, e ele curou a todos.
\VS{16}E ele lhes ordenava que não o tornassem conhecido;
\VS{17}Para que se cumprisse o que foi dito pelo profeta Isaías:
\VS{18}Eis aqui meu servo a quem escolhi, meu amado em quem minha alma se agrada; sobre ele porei o meu Espírito, e ele anunciará justiça às nações.
\VS{19}Ele não fará brigas, nem gritará; ninguém ouvirá pelas ruas a sua voz.
\VS{20}A cana rachada ele não quebrará, o pavio que fumega ele não apagará, até que conduza a justiça à vitória.
\VS{21}E em seu nome as nações terão esperança.
{\RQ{Isaías 42:1-4
}}
\VS{22}Então lhe trouxeram um endemoninhado cego e mudo; e ele o curou de tal maneira que o cego e mudo passou a falar e a ver.
\VS{23}E todas as multidões se admiravam e diziam: Não é este o Filho de Davi?
\VS{24}Mas quando os fariseus ouviam isso, diziam: Ele não expulsa os demônios, a não ser por Belzebu, o chefe dos demônios.
\VS{25}Porém Jesus, entendendo os pensamentos deles, disse-lhes: Todo reino dividido contra si mesmo é destruído; e toda cidade ou casa dividida contra si mesma não permanecerá.
\VS{26}Ora, se Satanás expulsa a Satanás, contra si mesmo está dividido; como, pois, permanecerá o seu reino?
\VS{27}E se eu expulso os demônios por Belzebu, então por quem vossos filhos
\FTNT{†}{{\FR 12:27 }
vossos filhos
i.e., vossos discípulos
} os expulsam? Portanto, eles mesmos serão vossos juízes.
\VS{28}Mas se eu pelo Espírito de Deus expulso os demônios, logo o Reino de Deus já chegou sobre vos.
\VS{29}Ou como pode alguém entrar na casa do valente, e saquear seus bens, sem primeiro amarrar ao valente? Depois disso saqueará sua casa.
\VS{30}Quem não é comigo é contra mim; e quem não ajunta comigo, espalha.
\VS{31}Por isso eu vos digo: todo pecado e blasfêmia serão perdoados aos seres humanos; mas a blasfêmia contra o Espírito não será perdoada aos seres humanos.
\VS{32}E qualquer um que falar palavra contra o Filho do homem lhe será perdoado; mas qualquer um que falar contra o Espírito Santo, não lhe será perdoado, nem na era presente, nem na futura.
\VS{33}Ou fazei a árvore boa, e seu fruto bom; ou fazei a árvore má, e seu fruto mau; pois pelo fruto se conhece a árvore.
\VS{34}Ninhada de víboras, como podeis vós falar boas coisas, sendo maus? Pois a boca fala do que o coração tem em abundância.
\VS{35}A pessoa boa tira coisas boas do bom tesouro do coração, e a pessoa má tira coisas más do tesouro mau.
\VS{36}Eu, porém, vos digo que de toda palavra imprudente
\FTNT{‡}{{\FR 12:36 }
imprudente
ou inútil
} que as pessoas falarem, dela prestarão contas no dia do juízo.
\VS{37}Porque por tuas palavras serás justificado,
\FTNT{§}{{\FR 12:37 }
justificado
ou absolvido
} e por tuas palavras serás condenado.
\VS{38}Então responderam uns dos escribas e dos fariseus, dizendo: Mestre, queremos ver de ti algum sinal.
\VS{39}Mas ele lhes deu a seguinte resposta: Uma geração má e adúltera pede sinal; mas não lhe será dado, exceto o sinal do profeta Jonas.
\VS{40}Porque assim como Jonas esteve três dias e três noites no ventre da baleia,
\FTNT{*}{{\FR 12:40 }
baleia
ou “grande animal marinho”, isto é, não necessariamente dos animais chamados atualmente de baleias.
} assim também o Filho do homem estará três dias e três noites no coração da terra.
\VS{41}Os de Nínive se levantarão no Juízo com esta geração, e a condenarão; porque se arrependeram com a pregação de Jonas. E eis aqui quem é maior que Jonas.
\VS{42}A rainha do sul se levantará no Juízo com esta geração, e a condenará; porque veio dos confins da terra para ouvir a sabedoria de Salomão. E eis aqui quem é maior que Salomão.
\VS{43}Quando o espírito imundo sai de alguém, anda por lugares secos buscando repouso, e não o acha.
\VS{44}Então diz: ‘Voltarei para a minha casa, de onde saí.’ E quando chega, a encontra desocupada, varrida, e arrumada.
\VS{45}Então vai, e toma consigo outros sete espíritos piores que ele; eles entram, e moram ali; e a última
{\ADD{condição}} de tal pessoa se torna pior que a primeira. Assim também acontecerá com esta geração má.
\VS{46}Enquanto ele ainda estava falando às multidões, eis que sua mãe e seus irmãos estavam fora, querendo falar com ele.
\VS{47}E alguém lhe disse: Eis que tua mãe e teus irmãos estão fora, querendo falar contigo.
\VS{48}Porém ele disse em resposta ao que o avisou: Quem é minha mãe? E quem são meus irmãos?
\VS{49}Então estendeu sua mão para os seus discípulos, e disse: Eis aqui minha mãe e meus irmãos;
\VS{50}Pois qualquer um que fizer a vontade do meu Pai, que está nos céus, esse é meu irmão, e irmã, e mãe.

\par }\Chap{13}{\PP \VerseOne{1}Naquele dia, Jesus saiu de casa e se sentou junto ao mar.
\VS{2}E ajuntaram-se perto dele tantas multidões, de maneira que ele entrou num barco e se sentou; e toda a multidão ficou na praia,
\VS{3}E ele lhes falou muitas coisas por parábolas. Ele disse: Eis que o semeador saiu a semear.
\VS{4}E enquanto semeava, caiu parte
{\ADD{das sementes}} junto ao caminho, e vieram as aves e a comeram.
\VS{5}E outra
{\ADD{parte}} caiu entre pedras, onde não havia muita terra, e logo nasceu, porque não tinha terra funda.
\VS{6}Mas quando o sol surgiu, queimou-se; e por não ter raiz, secou-se.
\VS{7}E outra
{\ADD{parte}} caiu entre espinhos, e os espinhos cresceram e a sufocaram.
\VS{8}E outra
{\ADD{parte}} caiu em boa terra, e rendeu fruto: um a cem, outro a sessenta, e outro a trinta.
\FTNT{*}{{\FR 13:8 }
cem, sessenta, trinta
isto é, produz cem, sessenta ou trinta vezes mais que o que foi semeado
}
\VS{9}Quem tem ouvidos para ouvir, ouça.
\VS{10}Então os discípulos se aproximaram, e lhe perguntaram: Por que falas a eles por parábolas?
\VS{11}E ele lhes respondeu: Porque a vós é dado saber os mistérios do Reino dos céus, mas a eles não é dado.
\VS{12}Pois a quem tem, lhe será dado, e terá em abundância; mas a quem não tem, até aquilo que tem lhe será tirado.
\VS{13}Por isso falo a eles por parábolas; porque vendo, não veem; e ouvindo, não ouvem, nem entendem.
\VS{14}Assim neles se cumpre a profecia de Isaías, que diz: De fato ouvireis, mas não entendereis; De fato vereis, mas não percebereis.
\VS{15}Porque o coração deste povo está insensível; Com seus ouvidos dificilmente ouvem, e seus olhos fecharam; A fim de não haver que seus olhos vejam, seus ouvidos ouçam, Seus corações entendam, e se arrependam, E eu os cure.
{\RQ{Isaías 6:9-10
}}
\VS{16}Mas benditos
\FTNT{†}{{\FR 13:16 }
benditos
ou “bem-aventurados”
} são os vossos olhos, porque veem; e os vossos ouvidos, porque ouvem.
\VS{17}Pois em verdade vos digo que muitos profetas e justos desejaram ver o que vós vedes, mas não viram; e
{\ADD{desejaram}} ouvir o que vós ouvis, mas não ouviram.
\VS{18}Portanto, ouvi vós a parábola do semeador:
\VS{19}Quando alguém ouve a palavra do Reino e não a entende, o maligno vem e arranca o que foi semeado em seu coração; este é o que foi semeado junto ao caminho.
\VS{20}E o que foi semeado entre as pedras é o que ouve a palavra, e logo a recebe com alegria,
\VS{21}mas não tem raiz em si mesmo. Em vez disso, dura um pouco, mas quando vem a aflição ou a perseguição pela palavra, logo tropeça na fé.
\VS{22}E o que foi semeado entre os espinhos é o que ouve a palavra, mas a ansiedade com o tempo presente e a sedução das riquezas sufocam a palavra, e fica sem dar fruto.
\VS{23}Mas o que foi semeado em boa terra, este é o que ouve e entende a palavra, e o que dá e produz fruto, um a cem, outro a sessenta, e outro a trinta.
\VS{24}E ele lhes declarou outra parábola, dizendo: O Reino dos céus é semelhante a um homem que semeia boa semente em seu campo,
\VS{25}Mas, enquanto as pessoas dormiam, o inimigo dele veio, semeou joio entre o trigo, e foi embora.
\VS{26}E, quando a erva cresceu e produziu fruto, então apareceu também o joio.
\VS{27}Então os servos do dono da propriedade chegaram, e lhe perguntaram: “Senhor, não semeaste boa semente no teu campo? De onde, pois, veio o joio?”
\VS{28}E ele lhes respondeu: “Um inimigo fez isto”. Em seguida, os servos lhe perguntaram: “Queres, pois, que vamos e o tiremos?”
\VS{29}Ele, porém, lhes respondeu: “Não, para não haver que, enquanto tirais o joio, arranqueis com ele também o trigo.
\VS{30}Deixai-os crescer ambos juntos até a colheita; e no tempo da colheita direi aos que colhem: ‘Recolhei primeiro o joio, e amarrai-o em molhos, para o queimarem; mas ao trigo ajuntai no meu celeiro’.”
\VS{31}Ele lhes propôs outra parábola: O Reino dos céus é semelhante a um grão de mostarda que alguém tomou e semeou no seu campo.
\VS{32}De fato, dentre todas as sementes, esta é a menor. Mas quando cresce, é a maior das hortaliças; e se torna
{\ADD{tamanha}} árvore, que as aves do céu vêm e se aninham em seus ramos.
\VS{33}Ele lhes disse outra parábola: O Reino dos céus é semelhante ao fermento que uma mulher tomou e misturou em três medidas de farinha,
\FTNT{‡}{{\FR 13:33 }
três medidas de farinha
equiv. Cerca de 40 litros (1 medida de farinha era aproximadamente 13 litros)
} até que tudo ficasse fermentado.
\VS{34}Tudo isto Jesus falou por parábolas às multidões. Sem parábolas ele não lhes falava,
\VS{35}para que se cumprisse o que foi falado pelo profeta, que disse: Abrirei a minha boca em parábolas; Pronunciarei coisas escondidas desde a fundação do mundo.
{\RQ{Salmos 78:2
}}
\VS{36}Então Jesus despediu as multidões, e foi para casa. Seus discípulos se aproximaram dele, e disseram: Explica-nos a parábola do joio do campo.
\VS{37}E ele lhes respondeu: O que semeia a boa semente é o Filho do homem.
\VS{38}E o campo é o mundo; e a boa semente, estes são os filhos do Reino; e o joio são os filhos do maligno.
\VS{39}E o inimigo, que o semeou, é o diabo; e a colheita é o fim da era;
\FTNT{§}{{\FR 13:39 }
fim da era
tradicionalmente “fim do mundo”, mas a palavra grega se refere ao período de tempo
} e os que colhem são os anjos.
\VS{40}Portanto, como o joio é colhido e queimado no fogo, assim também será no fim desta era.
\VS{41}O Filho do homem enviará seus anjos, e eles recolherão do seu Reino todas as causas do pecado,
\FTNT{*}{{\FR 13:41 }
causas do pecado
Tradicionalmente “escândalos”
} assim como os que praticam injustiça,
\VS{42}e os lançarão na fornalha de fogo. Ali haverá choro e ranger de dentes.
\VS{43}Então os justos brilharão como o sol, no Reino de seu Pai. Quem tem ouvidos para ouvir, ouça.
\VS{44}O Reino dos céus também é semelhante a um tesouro escondido num campo, que um homem, depois de achá-lo, escondeu. Então, em sua alegria, vai, vende tudo quanto tem, e compra aquele campo.
\VS{45}O Reino dos céus também é semelhante a um homem negociante, que buscava boas pérolas.
\VS{46}Quando este achou uma pérola de grande valor, foi, vendeu tudo quanto tinha, e a comprou.
\VS{47}O Reino dos céus também é semelhante a uma rede lançada ao mar, que colhe toda espécie
{\ADD{de peixes}} .
\VS{48}E quando está cheia,
{\ADD{os pescadores}} puxam-na à praia, sentam-se, e recolhem os bons em cestos, mas os ruins lançam fora.
\VS{49}Assim será ao fim da era; os anjos sairão, e separarão dentre os justos os maus,
\VS{50}e os lançarão na fornalha de fogo. Ali haverá choro e ranger de dentes.
\VS{51}E Jesus lhes perguntou: Entendestes todas estas coisas?Eles lhe responderam: Sim, Senhor.
\VS{52}E ele lhes disse: Portanto todo escriba que se tornou discípulo no Reino dos céus é semelhante a um chefe de casa, que do seu tesouro tira coisas novas e velhas.
\VS{53}E aconteceu que, quando Jesus acabou essas parábolas, retirou-se dali.
\VS{54}E vindo à sua terra, ensinava-os na sinagoga deles, de tal maneira que ficavam admirados, e diziam: De onde
{\ADD{vêm}} a este tal sabedoria, e os milagres?
\VS{55}Não é este o filho do carpinteiro? E não se chama sua mãe Maria, e seus irmãos Tiago, José, Simão, e Judas?
\VS{56}Não estão todas as suas irmãs conosco? Ora, de onde
{\ADD{vem}} a este tudo isto?
\VS{57}E ofenderam-se
\FTNT{†}{{\FR 13:57 }
ofenderam-se
Tradicionalmente: escandalizaram-se
} por causa dele. Mas Jesus lhes disse: Não há profeta sem honra, a não ser em sua terra, e em sua casa.
\VS{58}E não fez ali muitos milagres por causa da incredulidade deles.

\par }\Chap{14}{\PP \VerseOne{1}Naquele tempo Herodes, o tetrarca, ouviu relato a respeito de Jesus,
\VS{2}e disse aos seus servos: Este é João Batista; ele ressuscitou dos mortos, e por isso os milagres operam nele.
\VS{3}Porque Herodes havia prendido a João, acorrentado-o,
\FTNT{*}{{\FR 14:3 }
acorrentado
ou: amarrado
} e posto na prisão, por causa de Herodias, mulher do seu irmão Filipe;
\VS{4}pois João lhe dizia: Não te lícito que a tenhas.
\VS{5}{\ADD{Herodes}} queria matá-lo, mas tinha medo do povo, pois o consideravam profeta.
\VS{6}Porém, enquanto era celebrado o aniversário de Herodes, a filha de Herodias dançou no meio
{\ADD{das pessoas}} , e agradou a Herodes.
\VS{7}Por isso prometeu a ela dar tudo o que pedisse.
\VS{8}E ela, tendo sido induzida por sua mãe, disse: Dá-me aqui num prato a cabeça de João Batista.
\VS{9}E o rei se entristeceu; mas devido ao juramento, e aos que estavam presentes, ordenou que isso fosse concedido.
\VS{10}Então mandou degolarem João na prisão.
\VS{11}Sua cabeça foi trazida num prato, e dada à garota, e ela a levou à sua mãe.
\VS{12}E seus discípulos vieram, tomaram o corpo, e o enterraram; e foram avisar a Jesus.
\VS{13}Depois de Jesus ouvir, retirou-se dali num barco, a um lugar deserto, sozinho; mas assim que as multidões ouviram acerca disso, seguiram-no a pé das cidades.
\VS{14}Quando Jesus saiu, viu uma grande multidão. Ele se compadeceu deles, e curou dentre eles os enfermos.
\VS{15}E chegando o entardecer, os seus discípulos se aproximaram dele, e disseram: O lugar é deserto, e já é tarde.
\FTNT{†}{{\FR 14:15 }
já é tarde
lit. a hora [já] se passou
} Despede as multidões, para irem às aldeias, e comprarem para si de comer.
\VS{16}Mas Jesus lhes respondeu: Eles não precisam ir. Vós mesmos, dai-lhes de comer.
\VS{17}E eles lhe disseram: Nada temos aqui além de cinco pães e dois peixes.
\VS{18}Então disse: Trazei-os aqui a mim.
\VS{19}Ele mandou às multidões que se sentassem sobre a grama. Então tomou os cinco pães e os dois peixes, levantou os olhos ao céu, e
{\ADD{os}} abençoou. Em seguida partiu os pães, deu-os aos discípulos, e os discípulos às multidões.
\VS{20}E todos comeram, e se fartaram. E do que sobrou dos pedaços levantaram doze cestos cheios.
\VS{21}E os que comeram foram quase cinco mil homens, sem contar as mulheres e as crianças.
\VS{22}E logo Jesus mandou os seus discípulos entrarem no barco, e que fossem adiante dele para a outra margem, enquanto ele despedia as multidões.
\VS{23}Depois de despedir as multidões, subiu ao monte, à parte, para orar. Tendo chegado a noite, ele estava ali sozinho.
\VS{24}E o barco já estava no meio do mar, atormentado pelas ondas, porque o vento era contrário.
\VS{25}Mas à quarta vigília
\FTNT{‡}{{\FR 14:25 }
vigília
cada noite era dividida em quatro vigílias, de igual duração. Portanto a quarta vigília é, aproximadamente, das 3 horas da madrugada até o amanhecer
} da noite Jesus foi até eles, andando sobre o mar.
\VS{26}Quando os discípulos o viram andar sobre o mar, apavoraram-se, dizendo: É um fantasma! E gritaram de medo.
\VS{27}Mas Jesus logo lhes falou, dizendo: Tende coragem! Sou eu, não tenhais medo.
\VS{28}E Pedro lhe respondeu, dizendo: Senhor, se és tu, manda-me vir a ti sobre as águas.
\VS{29}E ele disse: Vem. Então Pedro desceu do barco e andou sobre as águas, para vir em direção a Jesus.
\VS{30}Mas quando viu o vento forte, teve medo; e começando a afundar, gritou: Senhor, salva-me!
\VS{31}Imediatamente Jesus estendeu a mão, segurou-o, e disse-lhe:
{\ADD{Homem}} de pouca fé, por que duvidaste?
\VS{32}E quando subiram no barco, o vento se aquietou.
\VS{33}Então os que estavam no barco vieram e o adoraram, dizendo: Verdadeiramente tu és o Filho de Deus.
\VS{34}E havendo passado para a outra margem, chegaram à terra de Genesaré.
\VS{35}E quando os homens daquele lugar o reconheceram, deram aviso por toda aquela região em redor, e lhe trouxeram todos os que estavam enfermos.
\VS{36}E rogavam-lhe que tão somente tocassem a borda
\FTNT{§}{{\FR 14:36 }
borda
ou: franja
} de sua roupa;
\FTNT{*}{{\FR 14:36 }
roupa
ou: capa
} e todos os que tocavam ficaram curados.

\par }\Chap{15}{\PP \VerseOne{1}Então
{\ADD{alguns}} escribas e fariseus de Jerusalém se aproximaram de Jesus, e perguntaram:
\VS{2}Por que os teus discípulos transgridem a tradição dos anciãos? Pois não lavam suas mãos quando comem pão.
\VS{3}Porém ele lhes respondeu: E vós, por que transgredis o mandamento de Deus por vossa tradição?
\VS{4}Pois Deus mandou, dizendo: Honra ao teu pai e à
{\ADD{tua}} mãe; e quem maldisser ao pai ou à mãe seja sentenciado à morte.
{\RQ{Êxodo 21:17, Levítico 20:9
}}
\VS{5}Mas vós dizeis: “Qualquer um que disser ao pai ou à mãe: ‘Todo o proveito que terias de mim é oferta exclusiva
{\ADD{para Deus}} ’, não
{\ADD{precisa}} honrar seu pai ou à sua mãe”.
\VS{6}E
{\ADD{assim}} invalidastes o mandamento de Deus por vossa tradição.
\VS{7}Hipócritas! Isaías bem profetizou sobre vós, dizendo:
\VS{8}Este povo com sua boca se aproxima de mim, e com os lábios me honra; mas o seu coração está longe de mim.
\VS{9}Em vão, porém, me veneram, ensinando doutrinas que são regras humanas.
{\RQ{Isaías 29:13
}}
\VS{10}Assim chamou a multidão para si, e disse-lhes: Ouvi e entendei.
\VS{11}Não é o que entra na boca que contamina o ser humano; mas sim o que sai da boca, isso contamina o ser humano.
\VS{12}Então os seus discípulos aproximaram-se dele, e lhe perguntaram: Tu sabes que os fariseus se ofenderam
\FTNT{*}{{\FR 15:12 }
Tradicionalmente: “se escandalizaram”
} quando ouviram essa palavra?
\VS{13}Mas ele respondeu: Toda planta que meu Pai celestial não plantou será arrancada pela raiz.
\VS{14}Deixai-os, são guias cegos de cegos. E se o cego guiar
{\ADD{outro}} cego, ambos cairão na cova.
\VS{15}E Pedro lhe disse: Explica-nos esta parábola.
\VS{16}Porém Jesus disse: Até vós ainda estais sem entender?
\VS{17}Não percebeis ainda que tudo o que entra na boca vai ao ventre, mas
{\ADD{depois}} é lançado na privada?
\VS{18}Porém as coisas que saem da boca procedem do coração; e elas contaminam o ser humano.
\VS{19}Pois do coração procedem maus pensamentos, mortes, adultérios, pecados sexuais, furtos, falsos testemunhos, blasfêmias.
\VS{20}Estas coisas são as que contaminam o ser humano; mas comer sem lavar as mãos não contamina o ser humano.
\VS{21}E, tendo Jesus partido dali, foi para as partes de Tiro e de Sidom.
\VS{22}E eis que uma mulher Cananeia, que tinha saído daquela região, clamou-lhe: Senhor, Filho de Davi, tem misericórdia de mim! Minha filha está miseravelmente endemoninhada.
\VS{23}Mas ele não lhe respondeu palavra. Então seus discípulos se aproximaram dele, e rogaram-lhe, dizendo: Manda-a embora, porque ela está gritando atrás de nós.
\VS{24}E ele respondeu: Não fui enviado para ninguém além das ovelhas perdidas
\FTNT{†}{{\FR 15:24 }
Não fui enviado para ninguém além das ovelhas perdidas
lit. Não fui enviado, a não ser para as ovelhas perdidas
} da casa de Israel.
\VS{25}Então ela veio e se prostrou diante dele, dizendo: Senhor, socorre-me.
\VS{26}Mas ele respondeu: Não é bom tomar o pão dos filhos e lançá-lo aos cachorrinhos.
\VS{27}Ela, porém, disse: Sim, Senhor. Porém os cachorrinhos também comem, das migalhas que caem da mesa dos seus senhores.
\VS{28}Então Jesus lhe respondeu: Ó mulher, grande
{\ADD{é}} a tua fé. A ti seja feito como tu queres. E desde aquela hora sua filha ficou curada.
\VS{29}E tendo Jesus partido dali, veio ao mar da Galileia. Ele subiu a um monte, e ali se sentou.
\VS{30}E vieram a ele muitas multidões, que tinham consigo mancos, cegos, mudos, aleijados, e muitos outros; e os lançaram aos pés de Jesus, e ele os curou.
\VS{31}Desta maneira, as multidões se maravilhavam quando viam os mudos falarem, os aleijados ficarem sãos, os mancos andarem, e os cegos verem; então glorificaram ao Deus de Israel.
\VS{32}Jesus chamou a si os seus discípulos, e disse: Estou compadecido com a multidão, porque já há três dias que estão comigo, e não têm o que comer. E não quero os deixar ir em jejum, para que não desmaiem no caminho.
\VS{33}E os seus discípulos lhe responderam: De onde conseguiremos tantos pães no deserto, para saciar tão grande multidão?
\VS{34}Jesus lhes perguntou: Quantos pães tendes? E eles disseram: Sete; e uns poucos peixinhos.
\VS{35}Então mandou as multidões que se sentassem pelo chão.
\VS{36}Tomou os sete pães e os peixes, deu graças e os partiu. Em seguida, ele os deu aos seus discípulos, e os discípulos à multidão.
\VS{37}E todos comeram e se saciaram; e levantaram dos pedaços que sobraram sete cestos cheios.
\VS{38}E foram os que comeram quatro mil homens, sem contar as mulheres e as crianças.
\VS{39}Depois de despedir as multidões,
{\ADD{Jesus}} entrou em um barco, e veio à região de Magdala.

\par }\Chap{16}{\PP \VerseOne{1}Então os fariseus e os saduceus se aproximaram dele e, a fim de tentá-lo, pediram-lhe que lhes mostrasse algum sinal do céu.
\VS{2}Mas ele lhes respondeu: Quando chega a tarde, dizeis: “
{\ADD{Haverá}} tempo bom, pois o céu está vermelho”.
\VS{3}E pela manhã: “Hoje
{\ADD{haverá}} tempestade, pois o céu está de um vermelho sombrio”. Hipócritas! Vós bem sabeis distinguir a aparência do céu, mas os sinais dos tempos não podeis?
\VS{4}Uma geração má e adúltera pede um sinal; mas nenhum sinal lhe será dado, a não ser o sinal do profeta Jonas. Então os deixou, e foi embora.
\VS{5}E quando os seus discípulos vieram para a outra margem, esqueceram-se de tomar pão.
\VS{6}E Jesus lhes disse: Ficai atentos, e tende cuidado com o fermento dos fariseus e saduceus.
\VS{7}E eles argumentaram entre si, dizendo: É porque não tomamos pão.
\VS{8}Jesus percebeu, e disse-lhes: Por que estais argumentando entre vós mesmos, ó
{\ADD{homens}} de pouca fé, que não tomastes pão?
\VS{9}Ainda não entendeis, nem vos lembrais dos cinco pães dos cinco mil, e quantos cestos levantastes?
\VS{10}Nem dos sete pães dos quatro mil, e quantos cestos levantastes?
\VS{11}Como não entendeis que não foi pelo pão que eu vos disse para tomardes cuidado com o fermento dos fariseus e saduceus?
\VS{12}Então entenderam que ele não havia dito que tomassem cuidado com o fermento de pão, mas sim, com a doutrina dos fariseus e saduceus.
\VS{13}E tendo Jesus vindo às partes da Cesareia de Filipe, perguntou aos seus discípulos: Quem as pessoas dizem que eu, o Filho do homem, sou?
\VS{14}E eles responderam: Alguns João Batista, outros Elias, e outros Jeremias ou algum dos profetas.
\VS{15}Ele lhes disse: E vós, quem dizeis que eu sou?
\VS{16}E Simão Pedro respondeu: Tu és o Cristo, o Filho do Deus vivo!
\VS{17}E Jesus lhe replicou: Bendito
\FTNT{*}{{\FR 16:17 }
Bendito
ou: bem-aventurado
} és tu, Simão, filho de Jonas;
\FTNT{†}{{\FR 16:17 }
filho de Jonas
lit. Bar-Jonas
} pois não foi carne e sangue que o revelou a ti, mas sim meu Pai, que
{\ADD{está}} nos céus.
\VS{18}E eu também te digo que tu és Pedro, e sobre esta pedra edificarei a minha igreja; e as portas do Xeol
\FTNT{‡}{{\FR 16:18 }
Xeol é o lugar dos mortos
} não prevalecerão contra ela.
\VS{19}E a ti darei as chaves do Reino dos céus; e tudo o que ligares na terra terá sido ligado
\FTNT{§}{{\FR 16:19 }
terá sido ligado
Ou: será ligado. No grego, as duas formas verbais são idênticas
} nos céus; e tudo o que desligares na terra terá sido desligado
\FTNT{*}{{\FR 16:19 }
ou: será desligado
} nos céus.
\VS{20}Então mandou aos seus discípulos que a ninguém dissessem que ele era Jesus, o Cristo.
\VS{21}Desde então Jesus começou a mostrar a seus discípulos que ele tinha que ir a Jerusalém, e sofrer muito pelos anciãos, pelos chefes dos sacerdotes, e pelos escribas, e ser morto, e ser ressuscitado ao terceiro dia.
\VS{22}E Pedro o tomou à parte, e começou a repreendê-lo, dizendo: Misericórdia de ti, Senhor! De maneira nenhuma isso te aconteça.
\VS{23}Mas ele se virou, e disse a Pedro: Para trás de mim, Satanás! Tu és um meio de tropeço, porque não compreendes as coisas de Deus, mas sim as humanas.
\VS{24}Então Jesus disse a seus discípulos: Se alguém quiser vir após mim, negue-se a si mesmo, tome sobre si a sua cruz, e siga-me.
\VS{25}Pois qualquer um que quiser salvar a sua vida a perderá; porém qualquer um que por causa de mim perder a sua vida,
{\ADD{este}} a achará.
\VS{26}Pois que proveito há para alguém, se ganhar o mundo todo, mas perder a sua alma? Ou que dará alguém em resgate da sua alma?
\VS{27}Pois o Filho do homem virá na glória do seu Pai com os seus anjos, e então recompensará a cada um segundo as suas obras.
\VS{28}Em verdade vos digo, que há alguns, dos que aqui estão, que não experimentarão a morte, até que vejam o Filho do homem vir em seu Reino.

\par }\Chap{17}{\PP \VerseOne{1}Seis dias depois, Jesus tomou consigo Pedro, Tiago, e seu irmão João, e os levou a sós a um monte alto.
\VS{2}Então transfigurou-se diante deles; seu rosto brilhou como o sol, e suas roupas se tornaram brancas como a luz.
\VS{3}E eis que lhes apareceram Moisés e Elias, falando com ele.
\VS{4}Pedro, então, disse a Jesus: Senhor, bom é para nós estarmos aqui. Se queres, façamos aqui três tendas:
\FTNT{*}{{\FR 17:4 }
tendas
ou: tabernáculos
} uma para ti, uma para Moisés, e uma para Elias.
\VS{5}Enquanto ele ainda estava falando, eis que uma nuvem brilhante os cobriu. E eis que uma voz da nuvem disse: Este é o meu Filho amado, em quem me agrado; a ele ouvi.
\VS{6}Quando os discípulos ouviram, caíram sobre seus rostos, e tiveram muito medo.
\VS{7}Jesus se aproximou deles, tocou-os, e disse: Levantai-vos, e não tenhais medo.
\VS{8}E quando eles levantaram seus olhos, não viram a ninguém, a não ser a Jesus somente.
\VS{9}E enquanto desciam do monte, Jesus lhes disse a seguinte ordem: Não conteis a visão a ninguém, até que o Filho do homem seja ressuscitado dos mortos.
\VS{10}E os seus discípulos lhe perguntaram: Por que, então, os escribas dizem que Elias tem que vir primeiro?
\VS{11}Jesus lhes respondeu: Em verdade Elias virá primeiro, e restaurará todas as coisas.
\VS{12}Digo-vos, porém, que Elias já veio, mas não o reconheceram. Em vez disso fizeram dele tudo o que quiseram. Assim também o Filho do homem sofrerá por meio deles.
\VS{13}Então os discípulos entenderam que ele lhes falara a respeito de João Batista.
\VS{14}E quando chegaram à multidão, veio a ele um homem, que se ajoelhou diante dele, e disse:
\VS{15}Senhor, tem misericórdia do meu filho, que é epilético,
\FTNT{†}{{\FR 17:15 }
epilético
tradicionalmente lunático
} e sofre muito mal; porque cai muitas vezes no fogo, e muitas vezes na água.
\VS{16}E eu o trouxe aos teus discípulos, mas não o puderam curar.
\VS{17}Jesus respondeu: Ó geração incrédula e perversa! Até quando estarei convosco? Até quando vos suportarei? Trazei-o a mim aqui.
\VS{18}E Jesus o repreendeu. Então o demônio saiu dele, e o menino sarou desde aquela hora.
\VS{19}Depois os discípulos se aproximaram de Jesus em particular, e perguntaram: Por que nós não o pudemos expulsar?
\VS{20}E Jesus lhes respondeu: Por causa da vossa incredulidade; pois em verdade vos digo, que se tivésseis fé como um grão de mostarda, diríeis a este monte: “Passa-te daqui para lá”, E ele passaria. E nada vos seria impossível.
\VS{21}Mas este tipo
{\ADD{de demônio}} não sai, a não ser por oração e jejum.
\VS{22}E enquanto eles estavam na Galileia, Jesus lhes disse: O Filho do homem será entregue em mãos de homens.
\VS{23}E o matarão, e ele será ressuscitado ao terceiro dia.E eles se entristeceram muito.
\VS{24}E quando entraram em Cafarnaum, os cobradores da taxa de duas dracmas
\FTNT{‡}{{\FR 17:24 }
taxa de duas dracmas
essa taxa era uma contribuição que os homens judeus faziam para financiar o Templo de Jerusalém. Uma dracma é equivalente a um denário, que era o pagamento de um dia de trabalho braçal
} vieram a Pedro, e perguntaram: Vosso mestre não paga as duas dracmas?
\VS{25}Ele respondeu: Sim. Quando ele entrou em casa, Jesus o antecipou, dizendo: Que te parece, Simão? De quem os reis da terra cobram tributos ou taxas? Dos seus filhos, ou dos outros?
\FTNT{§}{{\FR 17:25 }
outros: ou “estrangeiros”
}
\VS{26}Pedro lhe respondeu: Dos outros. Jesus lhe disse: Logo, os filhos são livres de pagar.
\VS{27}Mas para não os ofendermos, vai ao mar, e lança o anzol. Toma o primeiro peixe que subir, e quando lhe abrir a boca, acharás uma moeda de quatro dracmas. Toma-a, e dá a eles por mim e por ti.

\par }\Chap{18}{\PP \VerseOne{1}Naquela hora os discípulos se aproximaram de Jesus, e perguntaram: Ora, quem é o maior no Reino dos céus?
\VS{2}Então Jesus chamou a si uma criança, e a pôs no meio deles,
\VS{3}e disse: Em verdade vos digo, que se vós não converterdes, e fordes como crianças, de maneira nenhuma entrareis no Reino dos céus.
\VS{4}Assim, qualquer um que for humilde como esta criança, este é o maior no reino dos céus.
\VS{5}E qualquer um que receber a uma criança como esta em meu nome, recebe a mim.
\VS{6}Mas qualquer um que conduzir ao pecado
\FTNT{*}{{\FR 18:6 }
conduzir ao pecado
Tradicionalmente, escandalizar
} a um destes pequeninos que creem em mim, melhor lhe fora que uma grande pedra de moinho lhe fosse pendurada ao pescoço, e se afundasse no fundo do mar.
\VS{7}Ai do mundo por causa das tentações
\FTNT{†}{{\FR 18:7 }
Lit. tropeços. Tradicionalmente, escândalos
} do pecado! Pois é necessário que as tentações venham, mas ai daquela pessoa por quem a tentação vem!
\VS{8}Portanto, se a tua mão ou o teu pé te faz pecar,
\FTNT{‡}{{\FR 18:8 }
Lit. tropeçar. Também no v. 9
} corta-os, e lança-os de ti; melhor te é entrar manco ou aleijado na vida do que, tendo duas mãos ou dois pés, ser lançado no fogo eterno.
\VS{9}E se o teu olho te faz pecar, arranca-o, e lança-o de ti. Melhor te é entrar com um olho na vida do que, tendo dois olhos, ser lançado no inferno de fogo.
\VS{10}Olhai para que não desprezeis a algum destes pequeninos; porque eu vos digo que os seus anjos nos céus sempre veem a face do meu Pai, que
{\ADD{está}} nos céus.
\VS{11}Pois o Filho do homem veio para salvar o que havia se perdido.
\VS{12}Que vos parece? Se alguém tivesse cem ovelhas, e uma delas se desviasse, por acaso não iria ele pelos montes, deixando as noventa e nove, em busca da desviada?
\VS{13}E se acontecesse de achá-la, em verdade vos digo que ele se alegra mais daquela, do que das noventa e nove que se não desviaram.
\VS{14}Da mesma maneira, não é da vontade do vosso Pai, que
{\ADD{está}} nos céus, que um sequer destes pequeninos se perca.
\VS{15}Porém, se teu irmão pecar contra ti, vai repreendê-lo entre ti e ele só; se te ouvir, ganhaste o teu irmão.
\VS{16}Mas se não ouvir, toma ainda contigo um ou dois, para que toda palavra se confirme pela boca de duas ou três testemunhas.
{\RQ{Deuteronômio 19:15
}}
\VS{17}E se não lhes der ouvidos, comunica à igreja; e se também não der ouvidos à igreja, considera-o como gentio e publicano.
\VS{18}Em verdade vos digo que tudo o que vós ligardes na terra será ligado
\FTNT{§}{{\FR 18:18 }
será ligado
ou: terá sido ligado
} no céu; e tudo o que desligardes na terra será desligado
\FTNT{*}{{\FR 18:18 }
será desligado
ou: terá sido desligado
} no céu.
\VS{19}E digo-vos também que, se dois de vós concordarem na terra acerca de qualquer coisa que pedirem, isso lhes será feito por meu Pai, que
{\ADD{está}} nos céus.
\VS{20}Pois onde dois ou três estiverem reunidos em meu nome, ali eu estou no meio deles.
\VS{21}Então Pedro aproximou-se dele, e perguntou: Senhor, até quantas vezes meu irmão pecará contra mim, e eu lhe perdoarei? Até sete?
\VS{22}Jesus lhe respondeu: Eu não te digo até sete, mas sim até setenta vezes sete.
\VS{23}Por isso o Reino dos céus é comparável a um certo rei, que quis fazer acerto de contas com os seus servos.
\VS{24}E começando a fazer acerto de contas, foi-lhe apresentado um que lhe devia dez mil talentos.
\VS{25}Como ele não tinha com que pagar, o seu senhor mandou que ele, sua mulher, filhos, e tudo quanto tinha fossem vendidos para se fazer o pagamento.
\VS{26}Então aquele servo caiu e ficou prostrado diante dele, dizendo: “Senhor, tem paciência comigo, e tudo te pagarei”.
\VS{27}O senhor daquele servo compadeceu-se dele, então o soltou e lhe perdoou a dívida.
\VS{28}Todavia, depois daquele servo sair, achou um companheiro de serviço seu, que lhe devia cem denários;
\FTNT{†}{{\FR 18:28 }
denários
uma moeda de prata que valia um dia de trabalho braçal
} então o agarrou e o sufocou, dizendo: “Paga-me o que
{\ADD{me}} deves!”
\VS{29}Então o seu colega se prostrou diante dos seus pés, e lhe suplicou, dizendo: “Tem paciência comigo, e tudo te pagarei”.
\VS{30}Mas ele não quis. Em vez disso foi lançá-lo na prisão até que pagasse a dívida.
\VS{31}Quando os seus companheiros de serviço viram o que se passava, entristeceram-se muito. Então vieram denunciar ao seu senhor tudo o que havia se passado.
\VS{32}Assim o seu senhor o chamou, e lhe disse: “Servo mau! Toda aquela dívida te perdoei, porque me suplicaste.
\VS{33}Não tinhas tu a obrigação de ter tido misericórdia do sevo colega teu, assim como eu tive misericórdia de ti?”
\VS{34}E, enfurecido, o seu senhor o entregou aos torturadores até que pagasse tudo o que lhe devia.
\VS{35}Assim também meu Pai celestial vos fará, se não perdoardes de coração cada um ao seu irmão suas ofensas.

\par }\Chap{19}{\PP \VerseOne{1}E aconteceu que, quando Jesus acabou essas palavras, partiu da Galileia, e veio para a região da Judeia, além do Jordão.
\VS{2}E muitas multidões o seguiram, e ele os curou ali.
\VS{3}Então os fariseus se aproximaram dele e, provando-o, perguntaram-lhe: É lícito ao homem se divorciar da sua mulher por qualquer causa?
\VS{4}Porém ele lhes respondeu: Não tendes lido que aquele que
{\ADD{os}} fez no princípio, macho e fêmea os fez,
\VS{5}e disse: Portanto o homem deixará pai e mãe, e se unirá a sua mulher, e os dois serão uma única carne?
{\RQ{Gênesis 2:24
}}
\VS{6}Assim eles já não são mais dois, mas sim uma única carne; portanto, o que Deus juntou, o ser humano não separe.
\VS{7}Eles lhe disseram: Por que, pois, Moisés mandou
{\ADD{lhe}} dar carta de separação, e divorciar-se dela?
\VS{8}{\ADD{Jesus}} lhes disse: Por causa da dureza dos vossos corações Moisés vos permitiu divorciardes de vossas mulheres; mas no princípio não foi assim.
\VS{9}Porém eu vos digo que qualquer um que se divorciar de sua mulher, a não ser por causa de pecado sexual, e se casar com outra, adultera; e o que se casar com a divorciada
{\ADD{também}} adultera.
\VS{10}Os seus discípulos lhe disseram: Se assim é a condição do homem com a mulher, não convém se casar.
\VS{11}Porém ele lhes disse: Nem todos recebem esta palavra, a não ser aqueles a quem é dado;
\FTNT{*}{{\FR 19:11 }
quem é dado
i.e., quem é dado o dom ou capacidade de não ter relações sexuais
}
\VS{12}Pois há castrados que nasceram assim do ventre da mãe; e há castrados que foram castrados pelos homens; e há castrados que castraram a si mesmos por causa do Reino dos céus. Quem pode receber
{\ADD{isto}} ,receba.
\VS{13}Então lhe trouxeram crianças, para que pusesse as mãos sobre elas e orasse, mas os discípulos os repreendiam.
\FTNT{†}{{\FR 19:13 }
os repreendiam
i.e., repreendiam os que traziam as crianças, e não as crianças
}
\VS{14}Mas Jesus disse: Deixai as crianças, e não as impeçais de vir a mim, porque delas é o Reino dos céus.
\VS{15}Ele pôs as mãos sobre elas, e depois partiu-se dali.
\VS{16}E eis que alguém se aproximou, e perguntou-lhe: “Bom Mestre, que bem farei para eu ter a vida eterna?”
\VS{17}E ele lhe disse: “Por que me chamas bom? Ninguém há bom, a não ser um: Deus. Porém, se queres entrar na vida, guarda os mandamentos.”
\VS{18}Perguntou-lhe ele: “Quais?” E Jesus respondeu: “Não cometerás homicídio, não adulterarás, não furtarás, não darás falso testemunho;
\VS{19}honra ao teu pai e à
{\ADD{tua}} mãe; e amarás ao teu próximo como a ti mesmo.”
\VS{20}O rapaz lhe disse: “Tenho guardado tudo isso desde a minha juventude. Que me falta ainda?”
\VS{21}Disse-lhe Jesus: “Se queres ser completo, vai, vende o que tens, e dá aos pobres. Assim terás um tesouro no céu. Então vem, segue-me.”
\VS{22}Mas quando o rapaz ouviu essa palavra, foi embora triste, porque tinha muitos bens.
\VS{23}Jesus, então, disse aos seus discípulos: “Em verdade vos digo que dificilmente um rico entrará no reino dos céus.
\VS{24}Aliás, eu vos digo que é mais fácil um camelo passar pela abertura de uma agulha do que o rico entrar no reino de Deus.”
\VS{25}Quando os seus discípulos ouviram
{\ADD{isso}} , espantaram-se muito, e disseram: “Quem, pois, pode se salvar?”
\VS{26}Jesus olhou
{\ADD{para eles}} ,e lhes respondeu: Para os seres humanos, isto é impossível; mas para Deus tudo é possível.
\VS{27}Então Pedro se pôs a falar, e lhe perguntou: Eis que deixamos tudo, e te seguimos; o que, pois, conseguiremos ter?
\VS{28}E Jesus lhes disse: Em verdade vos digo que vós que me seguistes, na regeneração, quando o Filho do homem se sentar no trono de sua glória, vós também vos sentareis sobre doze tronos, para julgar as doze tribos de Israel.
\VS{29}E qualquer um que houver deixado casas, ou irmãos, ou irmãs, ou pai, ou mãe, ou mulher, ou filhos, ou terras por causa do meu nome, receberá cem vezes tanto,
\FTNT{‡}{{\FR 19:29 }
TR, RP: cem vezes tanto - N4: muitas vezes mais
} e herdará a vida eterna.
\VS{30}Porém muitos primeiros serão últimos; e últimos, primeiros.

\par }\Chap{20}{\PP \VerseOne{1}Pois o reino dos céus é semelhante a um homem, dono de propriedade, que saiu de madrugada para empregar trabalhadores para a sua vinha.
\VS{2}Ele entrou em acordo com os trabalhadores por um denário ao dia, e os mandou à sua vinha.
\VS{3}E quando saiu perto da hora terceira,
\FTNT{*}{{\FR 20:3 }
hora terceira
aproximadamente 9 horas da manhã
} viu outros que estavam desocupados na praça.
\VS{4}Então disse-lhes: “Ide vós também à vinha, e vos darei o que for justo”. E eles foram.
\VS{5}Saindo novamente perto da hora sexta e nona, fez o mesmo.
\FTNT{†}{{\FR 20:5 }
hora sexta e hora nona
aproximadamente meio-dia e 3 horas da tarde
}
\VS{6}E quando saiu perto da décima primeira hora,
\FTNT{‡}{{\FR 20:6 }
décima primeira hora
aproximadamente 5 horas da tarde
} achou outros que estavam desocupados, e lhes perguntou: “Por que estais aqui o dia todo desocupados?”
\VS{7}Eles lhe disseram: “Porque ninguém nos empregou”. Ele lhes respondeu: “Ide vós também à vinha, e recebereis o que for justo”.
\VS{8}E chegando o anoitecer, o senhor da vinha disse ao seu mordomo: “Chama aos trabalhadores, e paga-lhes o salário, começando dos últimos, até os primeiros”.
\VS{9}Então vieram os de cerca da hora décima primeira, e receberam um denário cada um.
\VS{10}Quando os primeiros vieram, pensavam que receberiam mais; porém eles também receberam um denário cada um.
\VS{11}Assim, ao receberem, murmuraram contra o chefe de casa,
\VS{12}dizendo: “Estes últimos trabalharam uma única hora, e tu os igualaste conosco, que suportamos a carga e o calor do dia”.
\VS{13}Ele, porém, respondeu a um deles: “Amigo, nada de errado estou fazendo contigo. Não concordaste tu comigo por um denário?
\VS{14}Toma o que é teu, e vai embora; e quero dar a este último tanto quanto a ti.
\VS{15}Acaso não me é lícito fazer do que é meu o que eu quiser? Ou o teu olho é mau, porque eu sou bom?”
\VS{16}Assim os últimos serão primeiros; e os primeiros, últimos; pois muitos são chamados, mas poucos escolhidos.
\VS{17}E enquanto Jesus subia a Jerusalém, tomou consigo os doze discípulos à parte no caminho, e lhes disse:
\VS{18}Eis que estamos subindo a Jerusalém, e o Filho do homem será entregue aos chefes dos sacerdotes e aos escribas, e o condenarão à morte.
\VS{19}E o entregarão aos gentios, para que dele escarneçam, e o açoitem, e crucifiquem; mas ao terceiro dia ressuscitará.
\VS{20}Então se aproximou dele a mãe dos filhos de Zebedeu, com os seus filhos. Ela o adorou para lhe pedir algo.
\VS{21}E ele lhe perguntou: O que queres? Ela lhe disse: Dá ordem
\FTNT{§}{{\FR 20:21 }
Dá ordem
lit. Dize
} para que estes meus dois filhos se sentem, um à tua direita e outro à tua esquerda, no teu Reino.
\VS{22}Porém Jesus respondeu: Não sabeis o que pedis. Podeis vós beber o cálice que eu beberei, e ser batizados com o batismo com que eu sou batizado? Eles lhe disseram: Podemos.
\VS{23}E ele lhes disse: De fato meu cálice bebereis, e com o batismo com que eu sou batizado sereis batizados; mas sentar-se à minha direita, e à minha esquerda, não me cabe concedê-lo, mas
{\ADD{será}} para os que por meu Pai está preparado.
\VS{24}E quando os dez ouviram
{\ADD{isso}} ,indignaram-se contra os dois irmãos.
\VS{25}Então Jesus os chamou a si, e disse: Vós bem sabeis que os chefes dos gentios os dominam, e os grandes usam de autoridade sobre eles.
\VS{26}Mas não será assim entre vós. Ao contrário, quem quiser se tornar grande entre vós seja o vosso assistente;
\VS{27}e quem quiser ser o primeiro entre vós seja o vosso servo;
\VS{28}assim como o Filho do homem não veio para ser servido, mas sim para servir, e para dar a sua vida em resgate por muitos.
\VS{29}Quando eles saíram de Jericó, uma grande multidão o seguiu.
\VS{30}E eis que dois cegos assentados junto ao caminho, ao ouvirem que Jesus passava, clamaram: Senhor, Filho de Davi, tem misericórdia de nós!
\VS{31}E a multidão os repreendia, para que se calassem, mas eles clamavam ainda mais: Senhor, Filho de Davi, tem misericórdia de nós!
\VS{32}Então Jesus parou, chamou-os, e perguntou: Que quereis que eu vos faça?
\VS{33}Eles lhe responderam: Senhor, que nossos olhos sejam abertos.
\VS{34}E Jesus, compadecido deles, tocou-lhes os olhos. E logo os olhos deles enxergaram, e o seguiram.

\par }\Chap{21}{\PP \VerseOne{1}E quando se aproximaram de Jerusalém, e chegaram a Betfagé, ao monte das Oliveiras, então Jesus mandou dois discípulos, dizendo-lhes:
\VS{2}Ide à aldeia em vossa frente, e logo achareis uma jumenta amarrada, e um jumentinho com ela; desamarra-a, e trazei-os a mim.
\VS{3}E se alguém vos disser algo, direis: “O Senhor precisa deles, mas logo os devolverá”.
\FTNT{*}{{\FR 21:3 }
devolverá
ou: “O Senhor precisa deles”. E logo ele os deixará ir. O texto grego é ambíguo
}
\VS{4}Ora, tudo isto aconteceu para que se cumprisse o que foi dito pelo profeta, que disse:
\VS{5}Dizei à filha de Sião: “Eis que o teu rei vem a ti, manso,
\FTNT{†}{{\FR 21:5 }
manso
ou humilde
} e sentado sobre um jumento; um jumentinho, filho de uma animal de carga”.
{\RQ{Zacarias 9:9
}}
\VS{6}Os discípulos foram, e fizeram como Jesus havia lhes mandado;
\VS{7}Então trouxeram a jumenta e o jumentinho, puseram as suas capas sobre eles, e fizeram
{\ADD{-no}} montar sobre elas.
\VS{8}E uma grande multidão estendia suas roupas pelo caminho, e outros cortavam ramos das árvores, e os espalhavam pelo caminho.
\VS{9}E as multidões que iam adiante, e as que seguiam, clamavam: Hosana ao Filho de Davi! Bendito o que vem no nome do Senhor! Hosana nas alturas!
\VS{10}Enquanto ele entrava em Jerusalém, toda a cidade se alvoroçou, perguntando: Quem é este?
\VS{11}E as multidões respondiam: Este é Jesus, o Profeta de Nazaré de Galileia.
\VS{12}Jesus entrou no Templo de Deus; então expulsou todos os que estavam vendendo e comprando no Templo, e virou as mesas dos cambistas e as cadeiras dos que vendiam pombas.
\VS{13}E disse-lhes: Está escrito: “Minha casa será chamada casa de oração”; mas vós a tornastes em covil de ladrões!
{\RQ{Isaías 56:7; Jeremias 7:11
}}
\VS{14}E cegos e mancos vieram a ele no Templo, e ele os curou.
\VS{15}Quando os chefes dos sacerdotes e os escribas viram as maravilhas que ele fazia, e as crianças gritando no Templo: “Hosana ao Filho de Davi!”, eles ficaram indignados.
\VS{16}E perguntaram-lhe: Ouves o que estas
{\ADD{crianças}} dizem? E Jesus lhes respondeu: Sim. Nunca lestes: “Da boca das crianças e dos bebês
\FTNT{‡}{{\FR 21:16 }
bebês
i.e. dos que mamam
} providenciaste o louvor?”
{\RQ{Salmos 8:2
}}
\VS{17}Então ele os deixou, e saiu da cidade para Betânia, e ali passou a noite.
\VS{18}E pela manhã, enquanto voltava para a cidade, teve fome.
\VS{19}Quando ele viu uma figueira perto do caminho, veio a ela, mas nada nela achou, a não ser somente folhas. E disse-lhe: Nunca de ti nasça fruto, jamais!E imediatamente a figueira se secou.
\VS{20}Os discípulos viram, e ficaram maravilhados, dizendo: Como a figueira se secou de imediato?
\VS{21}Porém Jesus lhes respondeu: Em verdade vos digo: se tiverdes fé, e não duvidardes, vós não somente fareis isto à figueira, mas até se disserdes a este monte: “Levanta-te, e lança-te no mar”, isso se fará.
\VS{22}E tudo o que pedirdes em oração, crendo, recebereis.
\VS{23}Depois de entrar no templo, quando ele estava ensinando, os chefes dos sacerdotes e os anciãos do povo se aproximaram dele, perguntando: Com que autoridade fazes isto? E quem te deu esta autoridade?
\VS{24}Jesus lhes respondeu: Eu também vos farei uma pergunta. Se vós a responderdes a mim, também eu vos responderei com que autoridade faço isto.
\VS{25}De onde era o batismo de João? Do céu, ou dos seres humanos?E eles pensaram entre si mesmos, dizendo: Se dissermos: “Do céu”, ele nos dirá: “Por que, então, não crestes nele?
\VS{26}Mas se dissermos: “Dos seres humanos”, temos medo da multidão, pois todos consideram João como profeta.
\VS{27}Então responderam a Jesus: Não sabemos. E ele lhes disse: Nem eu vos digo com que autoridade faço isto.
\VS{28}Mas que vos parece? Um homem tinha dois filhos. Aproximando-se do primeiro, disse: “Filho, vai hoje trabalhar na minha vinha.”
\VS{29}Porém ele respondeu: “Não quero”; mas depois se arrependeu, e foi.
\VS{30}E, aproximando-se do segundo, disse da mesma maneira. E ele respondeu: “Eu
{\ADD{vou}} , senhor”, mas não foi.
\VS{31}Qual dos dois fez a vontade do pai? Eles lhe responderam: O primeiro. Jesus lhes disse: Em verdade vos digo que os publicanos e as prostitutas estão indo adiante de vós ao Reino de Deus.
\VS{32}Pois João veio a vós mesmos no caminho de justiça, mas não crestes nele; enquanto que os publicanos e as prostitutas nele creram. Vós, porém, mesmo tendo visto
{\ADD{isto}} , nem assim vos arrependestes, a fim de nele crer.
\VS{33}Ouvi outra parábola. Havia um homem, dono de uma propriedade. Ele plantou uma vinha, cercou-a, fundou nela uma prensa de uvas, e construiu uma torre. Depois a arrendou a uns lavradores, e partiu-se para um lugar distante.
\VS{34}Quando chegou o tempo dos frutos, enviou seus servos aos lavradores, para receberem os frutos que a ele pertenciam.
\VS{35}Mas os lavradores tomaram os seus servos, e feriram um, mataram outro, e apedrejaram outro.
\VS{36}Outra vez enviou outros servos, em maior número que os primeiros, mas fizeram-lhes o mesmo.
\VS{37}E por último lhes enviou o seu filho, dizendo: “Respeitarão ao meu filho”.
\VS{38}Mas quando os lavradores viram o filho, disseram entre si: “Este é o herdeiro. Venhamos matá-lo, e tomemos a sua herança”.
\VS{39}Então o agarraram, lançaram-no para fora da vinha, e o mataram.
\VS{40}Ora, quando o senhor da vinha chegar, o que fará com aqueles lavradores?
\VS{41}Eles lhe responderam: Aos maus dará uma morte má, e arrendará a vinha a outros lavradores, que lhe deem os frutos em seus tempos
{\ADD{de colheita}} .
\VS{42}Jesus lhes disse: Nunca lestes nas Escrituras: “A pedra que os construtores rejeitaram, essa se tornou cabeça da esquina. Isto foi feito pelo Senhor, e é maravilhoso aos nossos olhos”?
\FTNT{§}{{\FR 21:42 }
cabeça de esquina
i.e. a pedra angular, a principal da construção
}
{\RQ{Salmos 118:22,23
}}
\VS{43}Portanto eu vos digo que o reino de Deus será tirado de vós, e será dado a um povo que produza os frutos dele.
\FTNT{*}{{\FR 21:43 }
dele
i.e. do reino
}
\VS{44}E quem cair sobre esta pedra será quebrado; mas sobre quem ela cair, ela o tornará em pó.
\VS{45}Quando os chefes dos sacerdotes e os fariseus ouviram estas suas parábolas, entenderam que
{\ADD{Jesus}} estava falando deles.
\VS{46}E procuravam prendê-lo, mas temeram as multidões, pois elas o consideravam profeta.

\par }\Chap{22}{\PP \VerseOne{1}Então Jesus voltou a lhes falar por parábolas, dizendo:
\VS{2}O reino dos céus é semelhante a um rei que fez uma festa de casamento para o seu filho;
\VS{3}e mandou a seus servos que chamassem os convidados para a festa de casamento, mas não quiseram vir.
\VS{4}Outra vez ele mandou outros servos, dizendo: “Dizei aos convidados: ‘Eis que já preparei meu jantar: meus bois e animais cevados já foram mortos, e tudo está pronto. Vinde à festa de casamento’ ”.
\VS{5}Porém eles não deram importância e foram embora, um ao seu campo, e outro ao seu comércio;
\VS{6}e outros agarraram os servos dele, e os humilharam e os mataram.
\VS{7}Quando o rei Então enviou os seus exércitos, destruiu aqueles homicidas, e incendiou a cidade deles.
\VS{8}Em seguida, disse aos seus servos: “Certamente a festa de casamento está pronta, porém os convidados não eram dignos.
\VS{9}Ide, pois, às saídas dos caminhos, e convidai à festa de casamento tantos quantos achardes.
\VS{10}Aqueles servos saíram pelos caminhos, e ajuntaram todos quantos acharam, tanto maus como bons; e a festa de casamento se encheu de convidados.
\VS{11}Mas quando o rei entrou para ver os convidados, percebeu ali um homem que não estava vestido com roupa adequada para a festa de casamento.
\VS{12}Então lhe perguntou: “Amigo, como entraste aqui sem ter roupa para a festa?” E ele emudeceu.
\VS{13}Então o rei disse aos servos: “Amarrai-o nos pés e nas mãos, tomai-o, e lançai-o nas trevas de fora. Ali haverá pranto e o ranger de dentes”.
\VS{14}Pois muitos são chamados, porém poucos escolhidos.
\VS{15}Então os fariseus foram embora, e se reuniram para tramar como o apanhariam em cilada por algo que dissesse.
\FTNT{*}{{\FR 22:15 }
por algo que dissesse
lit. em [alguma] palavra
}
\VS{16}Depois lhe enviaram seus discípulos, juntamente com os apoiadores de Herodes, e perguntaram: Mestre, bem sabemos que tu és verdadeiro, e que com verdade ensinas o caminho de Deus, e que não te importas com a opinião de ninguém, porque não dás atenção à aparência humana.
\VS{17}Dize-nos, pois, o que te parece: é lícito dar tributo a César, ou não?
\VS{18}Mas Jesus, entendendo a sua malícia, disse: Por que me tentais,
\FTNT{†}{{\FR 22:18 }
tentais
ou: testais
} hipócritas?
\VS{19}Mostrai-me a moeda do tributo.E eles lhe trouxeram um denário.
\VS{20}E ele lhes perguntou: De quem é esta imagem, e a inscrição?
\VS{21}Eles lhe responderam: De César. Então ele lhes disse: Dai, pois, a César o que é de César, e a Deus o que é de Deus.
\VS{22}Quando ouviram isso, eles ficaram admirados; então o deixaram e se retiraram.
\VS{23}Naquele mesmo dia chegaram a ele os saduceus, que dizem não haver ressurreição, e perguntaram-lhe,
\VS{24}dizendo: Mestre, Moisés disse: Se um homem morrer sem ter filhos, seu irmão se casará com sua mulher, e gerará descendência
\FTNT{‡}{{\FR 22:24 }
gerará descendência
lit. levantará semente
} ao seu irmão.
{\RQ{Deuteronômio 25:5
}}
\VS{25}Ora, havia entre nós sete irmãos. O primeiro se casou, e depois morreu; e sem ter tido filhos,
\FTNT{§}{{\FR 22:25 }
filhos
lit. semente
} deixou sua mulher ao seu irmão.
\VS{26}E da mesma maneira também foi com o segundo, o terceiro, até os sete.
\VS{27}Por último, depois de todos, a mulher também morreu.
\VS{28}Assim, na ressurreição, a mulher será de qual dos sete? Pois todos a tiveram.
\VS{29}Jesus, porém, lhes respondeu: Errais, por não conhecerdes as Escrituras, nem o poder de Deus.
\VS{30}Porque na ressurreição, nem se tomam, nem se dão em casamento; mas são como os anjos de Deus no céu.
\VS{31}E sobre a ressurreição dos mortos, não lestes o que Deus vos falou:
\VS{32}Eu sou o Deus de Abraão, o Deus de Isaque, e o Deus de Jacó?
{\RQ{Êxodo 3:6
}} Deus não é Deus dos mortos, mas sim dos vivos!
\VS{33}Quando as multidões ouviram
{\ADD{isto}} ,ficaram admiradas de sua doutrina.
\VS{34}E os fariseus, ao ouvirem que ele havia feito os saduceus se calarem, reuniram-se.
\VS{35}E um deles, especialista da Lei, tentando-o, perguntou-lhe:
\VS{36}Mestre, qual é o grande mandamento na Lei?
\VS{37}E Jesus lhe respondeu: Amarás ao Senhor teu Deus com todo o teu coração, com toda a tua alma, e com todo o teu entendimento:
{\RQ{Deuteronômio 6:5
}}
\VS{38}este é o primeiro e grande mandamento.
\VS{39}E o segundo, semelhante a este,
{\ADD{é}} : Amarás o teu próximo como a ti mesmo.
{\RQ{Levítico 19:18
}}
\VS{40}Destes dois mandamentos dependem toda a Lei e os Profetas.
\VS{41}E, estando os fariseus reunidos, Jesus lhes perguntou,
\VS{42}dizendo: Que pensais vós acerca do Cristo? De quem ele é filho? Eles lhe responderam: De Davi.
\VS{43}{\ADD{Jesus}} lhes disse: Como, pois, Davi, em espírito, o chama Senhor, dizendo:
\VS{44}Disse o Senhor a meu Senhor: “Senta-te à minha direita, até que eu ponha os teus inimigos como estrado de teus pés”.
{\RQ{Salmos 110:1
}}
\VS{45}Ora, se Davi o chama Senhor, como é seu filho?
\VS{46}E ninguém podia lhe responder palavra; nem ninguém ousou desde aquele dia a mais lhe perguntar.

\par }\Chap{23}{\PP \VerseOne{1}Então Jesus falou às multidões e aos seus discípulos,
\VS{2}dizendo: Os escribas e os fariseus se sentam sobre o assento de Moisés.
\VS{3}Portanto, tudo o que eles vos disserem que guardeis, guardai, e fazei. Mas não façais segundo as suas obras, porque eles dizem e não fazem.
\VS{4}Pois eles amarram cargas pesadas e difíceis de levar, e as põem sobre os ombros das pessoas; porém eles nem sequer com o seu dedo as querem mover.
\VS{5}E fazem todas as suas obras a fim de serem vistos pelas pessoas: por isso alargam seus filactérios,
\FTNT{*}{{\FR 23:5 }
filactérios
pequenas bolsas com trechos do Antigo Testamento que os judeus usavam nos braços e na testa, especialmente quando oravam (Deuteronômio 6:8
} ) e fazem compridas as franjas de suas roupas.
\VS{6}Eles amam os primeiros assentos nas ceias, as primeiras cadeiras nas sinagogas,
\VS{7}as saudações nas praças, e serem chamados: “Rabi, Rabi” pelas pessoas.
\VS{8}Mas vós, não sejais chamados Rabi, porque o vosso Mestre é um: o Cristo; e todos vós sois irmãos.
\VS{9}E não chameis a ninguém na terra vosso pai; porque o vosso Pai é um: aquele que está nos céus.
\VS{10}Nem sejais chamados mestres; porque o vosso mestre é um: o Cristo.
\VS{11}Porém o maior de vós será vosso servo.
\VS{12}E o que a si mesmo se exaltar será humilhado; e o que a si mesmo se humilhar será exaltado.
\VS{13}Mas ai de vós, escribas e fariseus, hipócritas! Porque fechais o Reino dos céus em frente das pessoas; pois nem vós entrais, nem permitis a entrada do que estão para entrar.
\VS{14}Ai de vós, escribas e fariseus, hipócritas! Porque devorais as casas das viúvas, e
{\ADD{isso}} com pretexto de longas orações; por isso recebereis mais grave condenação.
\VS{15}Ai de vós, escribas e fariseus, hipócritas! Porque rodeais o mar e a terra para fazerdes um prosélito;
\FTNT{†}{{\FR 23:15 }
prosélito
pessoa não-israelita que se convertia ao judaísmo
} e quando é feito, vós o tornais filho do inferno duas vezes mais que a vós.
\VS{16}Ai de vós, guias cegos, que dizeis: “Qualquer um que jurar pelo templo, nada é; mas qualquer um que jurar pelo ouro do templo, devedor é”.
\VS{17}Tolos e cegos! Pois qual é maior: o ouro, ou o templo que santifica o ouro?
\VS{18}Também
{\ADD{dizeis}} : “Qualquer um que jurar pelo altar, nada é; mas quem jurar pela oferta que
{\ADD{está}} sobre ele, devedor é”.
\VS{19}Tolos e cegos! Pois qual é maior: a oferta, ou o altar que santifica a oferta?
\VS{20}Portanto, quem jurar pelo altar, jura por ele, e por tudo o que está sobre ele.
\VS{21}E quem jurar pelo templo, jura por ele, e por aquele que nele habita.
\VS{22}E quem jurar pelo Céu, jura pelo trono de Deus, e por aquele que sobre ele está sentado.
\VS{23}Ai de vós, escribas e fariseus, hipócritas! Porque dais o dízimo da hortelã, do endro, e do cominho, e desprezais o que é mais importante da Lei: a justiça, a misericórdia, e a fidelidade;
\FTNT{‡}{{\FR 23:23 }
fidelidade
ou: fé
} estas coisas devem ser feitas, sem se desprezar as outras.
\VS{24}Guias cegos, que coais um mosquito, e engolis um camelo!
\VS{25}Ai de vós, escribas e fariseus, hipócritas! Porque limpais o exterior do copo ou do prato, mas por dentro estão cheios de extorsão e cobiça.
\FTNT{§}{{\FR 23:25 }
cobiça
ou: falta de moderação
}
\VS{26}Fariseu cego! Limpa primeiro o interior do copo e do prato, para que também o exterior deles fique limpo.
\VS{27}Ai de vós, escribas e fariseus, hipócritas! Porque sois semelhantes aos sepulcros caiados, que por fora realmente parecem belos, mas por dentro estão cheios de ossos de cadáveres, e de toda imundícia.
\VS{28}Assim também vós, por fora, realmente pareceis justos às pessoas, mas por dentro estais cheios de hipocrisia e de injustiça.
\VS{29}Ai de vós, escribas e fariseus, hipócritas! Porque edificais os sepulcros dos profetas, adornais os monumentos dos justos,
\VS{30}e dizeis: “Se estivéssemos nos dias dos nossos pais, nunca teríamos sido cúmplices deles quando derramaram o sangue
\FTNT{*}{{\FR 23:30 }
quando derramaram o sangue
lit. no sangue
} dos profetas”.
\VS{31}Assim vós mesmos dais testemunho de que sois filhos dos que mataram os profetas.
\VS{32}Completai, pois, a medida de vossos pais.
\FTNT{†}{{\FR 23:32 }
Completai, pois, a medida de vossos pais
i.e. Completai de fazer o que vossos ancestrais começaram, (matando os justos que anunciavam a palavra de Deus, e em especial, matando Cristo)
}
\VS{33}Serpentes, ninhada de víboras! Como escapareis da condenação do inferno?
\VS{34}Por isso, eis que eu vos envio profetas, sábios, e escribas; a uns deles matareis e crucificareis, e a
{\ADD{outros}} deles açoitareis em vossas sinagogas, e perseguireis de cidade em cidade;
\VS{35}para que venha sobre vós todo o sangue justo que foi derramado sobre a terra, desde o sangue do justo Abel, até o sangue de Zacarias, filho de Baraquias, ao qual matastes entre o templo e o altar.
\VS{36}Em verdade vos digo que tudo isto virá sobre esta geração.
\VS{37}Jerusalém, Jerusalém, que matas os profetas, e apedrejas os que te são enviados! Quantas vezes eu quis ajuntar os teus filhos, como a galinha ajunta os seus pintos debaixo das asas; porém não quisestes!
\VS{38}Eis que vossa casa vos será deixada desolada.
\FTNT{‡}{{\FR 23:38 }
desolada
ou: deserta
}
\VS{39}Pois eu vos digo que a partir de agora não me vereis, até que digais: “Bendito aquele que vem no nome do Senhor”.
{\RQ{Salmos 118:26
}}

\par }\Chap{24}{\PP \VerseOne{1}Jesus saiu do templo, e se foi. Então seus discípulos se aproximaram dele para lhe mostrarem os edifícios do complexo do templo.
\VS{2}Mas Jesus lhes disse: Não vedes tudo isto? Em verdade vos digo, que não será deixada aqui pedra sobre pedra, que não seja derrubada.
\VS{3}E, depois de se assentar no monte das Oliveiras, os discípulos se aproximaram dele reservadamente, perguntando: Dize-nos, quando serão estas coisas, e que sinal haverá da tua vinda, e do fim da era?
\FTNT{*}{{\FR 24:3 }
fim da era
ou: “fim dos tempos”, “fim do mundo”
}
\VS{4}E Jesus lhes respondeu: Permanecei atentos, para que ninguém vos engane.
\VS{5}Porque muitos virão em meu nome, dizendo: “Eu sou o Cristo”, e enganarão a muitos.
\VS{6}E ouvireis de guerras, e de rumores de guerras. Olhai que não vos espanteis; porque é necessário, que tudo
{\ADD{isto}} aconteça, mas ainda não é o fim.
\VS{7}Pois se levantará nação contra nação, e reino contra reino; e haverá fomes, pestilências, e terremotos em diversos lugares.
\VS{8}Mas todas estas coisas são o começo das dores.
\VS{9}Então vos entregarão para serdes afligidos, e vos matarão; e sereis odiados por todas as nações, por causa do meu nome.
\VS{10}E muitos tropeçarão na fé;
\FTNT{†}{{\FR 24:10 }
tropeçarão na fé
Tradicionalmente, se escandalizarão
} e trairão uns aos outros, e uns aos outros se odiarão.
\VS{11}E muitos falsos profetas se levantarão, e enganarão a muitos.
\VS{12}E, por se multiplicar a injustiça, o amor de muitos se esfriará.
\VS{13}Mas o que perseverar até o fim, esse será salvo.
\VS{14}E este Evangelho do Reino será pregado em todo o mundo, como testemunho a todas as nações, e então virá o fim.
\VS{15}Portanto, quando virdes que a abominação da desolação, dita pelo profeta Daniel, está no lugar santo, (quem lê, entenda),
\VS{16}então os que estiveram na Judeia fujam para os montes;
\VS{17}o que estiver no sobre o telhado não desça para tirar alguma coisa de sua casa;
\VS{18}e o que estiver no campo não volte atrás para tomar as suas roupas.
\VS{19}Mas ai das grávidas e das que amamentarem naqueles dias!
\VS{20}Orai, porém, para que a vossa fuga não aconteça no inverno, nem no sábado.
\VS{21}Pois haverá então grande aflição, como nunca houve desde o princípio do mundo até agora, nem jamais haverá.
\FTNT{‡}{{\FR 24:21 }
grande aflição: ou “grande tribulação”
}
\VS{22}E se aqueles dias não fossem encurtados, ninguém
\FTNT{§}{{\FR 24:22 }
ninguém
lit. “nenhuma carne”
} se salvaria;
\FTNT{*}{{\FR 24:22 }
se salvaria
i.e. sobreviveria
} mas por causa dos escolhidos, aqueles dias serão encurtados.
\VS{23}Então, se alguém vos disser: “Olha o Cristo aqui”, ou “
{\ADD{Olha ele}} ali”, não creiais,
\VS{24}pois se levantarão falsos cristos e falsos profetas; e farão tão grandes sinais e prodígios que, se fosse possível, enganariam até os escolhidos.
\VS{25}Eis que eu tenho vos dito com antecedência.
\VS{26}Portanto, se vos disserem: “Eis que ele está no deserto”, não saiais; “Eis que ele está em um recinto”, não creiais.
\VS{27}Porque, assim como o relâmpago, que sai do oriente, e aparece até o ocidente, assim também será a vinda do Filho do homem.
\VS{28}Pois onde estiver o cadáver, ali se ajuntarão os abutres.
\VS{29}E logo depois da aflição
\FTNT{†}{{\FR 24:29 }
aflição
ou “tribulação”
} daqueles dias, o sol se escurecerá, a lua não dará o seu brilho, as estrelas cairão do céu, e as forças dos céus se estremecerão.
\VS{30}Então aparecerá no céu o sinal do Filho do homem. Naquela hora todas as tribos da terra lamentarão, e verão ao Filho do homem, que vem sobre as nuvens do céu, com poder e grande glória.
\VS{31}E enviará os seus anjos com grande som de trombeta, e ajuntarão os seus escolhidos desde os quatro ventos, de uma extremidade à outra dos céus.
\VS{32}Aprendei a parábola da figueira: “Quando os seus ramos já ficam verdes, e as folhas brotam, sabeis que o verão
{\ADD{está}} perto”.
\VS{33}Assim também vós, quando virdes todas estas coisas, sabei que já está perto, às portas.
\VS{34}Em verdade vos digo que esta geração não passará, até que todas estas coisas aconteçam.
\VS{35}O céu e a terra passarão, mas minhas palavras de maneira nenhuma passarão.
\VS{36}Porém daquele dia e hora, ninguém sabe, nem os anjos do céu, a não ser meu Pai somente.
\VS{37}Assim como foram os dias de Noé, assim também será a vinda do Filho do homem.
\VS{38}Pois, assim como nos dias antes do dilúvio comiam, bebiam, casavam, e davam-se em casamento, até o dia em que Noé entrou na arca;
\VS{39}e não sabiam, até que veio o dilúvio, e levou todos, assim também será a vinda do Filho do homem.
\VS{40}Naquela hora dois estarão no campo; um será tomado, e o outro será deixado.
\VS{41}Duas estarão moendo em um moinho; uma será tomada, e a outra será deixada.
\VS{42}Vigiai, pois, porque não sabeis em que hora o vosso Senhor virá.
\VS{43}Porém sabei isto: se o dono de casa soubesse a que hora
\FTNT{‡}{{\FR 24:43 }
hora
lit. vigília – a noite era dividida em quatro vigílias, cada uma com cerca de três horas de duração
} da noite o ladrão viria, vigiaria, e não deixaria invadir a sua casa.
\VS{44}Portanto também vós estai prontos, porque o Filho do homem virá na hora que não esperais.
\VS{45}Pois quem é o servo fiel e prudente, ao qual o seu senhor pôs sobre os seus trabalhadores, para
{\ADD{lhes}} dar alimento no tempo devido?
\VS{46}Feliz
\FTNT{§}{{\FR 24:46 }
Feliz
ou: “bem-aventurado”
} será aquele servo a quem, quando o seu senhor vier, achar fazendo assim.
\VS{47}Em verdade vos digo que ele o porá sobre todos os seus bens.
\VS{48}Porém se aquele servo mau disser em seu coração: “Meu senhor está demorando a chegar”,
\VS{49}e começar a espancar os
{\ADD{seus}} companheiros de serviço, e a comer, e a beber com os beberrões,
\VS{50}o senhor daquele servo chegará num dia que ele não espera, e numa hora que ele não sabe,
\VS{51}e o despedaçará, e porá sua parte com os hipócritas; ali haverá pranto e ranger de dentes.

\par }\Chap{25}{\PP \VerseOne{1}Então o Reino dos céus será semelhante a dez virgens, que tomaram suas lâmpadas, e saíram ao encontro do noivo.
\VS{2}E cinco delas eram prudentes, e cinco tolas.
\VS{3}As tolas, quando tomaram as suas lâmpadas, não tomaram azeite consigo.
\VS{4}Mas as prudentes tomaram azeite nos seus frascos, com as suas lâmpadas.
\VS{5}O noivo demorou, por isso todas cochilaram e adormeceram.
\VS{6}Mas à meia-noite houve um grito: “Eis que vem o noivo! Ide ao seu encontro!”.
\VS{7}Então todas aquelas virgens se levantaram, e prepararam suas lâmpadas.
\VS{8}E as tolas disseram às prudentes: “Dai-nos do vosso azeite, porque as nossas lâmpadas estão se apagando”.
\VS{9}Mas as prudentes responderam: “
{\ADD{Não}} ,para que não falte a nós e a vós; em vez disso, ide aos vendedores, e comprai para vós mesmas”.
\VS{10}Enquanto elas foram comprar, veio o noivo. As que estavam preparadas entraram com ele à festa do casamento, e fechou-se a porta.
\VS{11}Depois vieram também as outras virgens, dizendo: “Senhor, Senhor, abre-nos!”
\VS{12}Mas ele respondeu: “Em verdade vos digo que não vos conheço”.
\VS{13}Portanto, vigiai, porque não sabeis o dia nem a hora em que o Filho do homem virá.
\VS{14}Pois
{\ADD{é}} como um homem, que partindo para fora do país, chamou seus servos, e lhes entregou os seus bens.
\VS{15}E a um deu cinco talentos, a outro dois, e ao terceiro um, a cada um conforme a sua habilidade, e logo depois partiu em viagem.
\VS{16}Em seguida, o que havia recebido cinco talentos foi fazer negócios com eles, e obteve outros cinco talentos.
\VS{17}E, semelhantemente, o que
{\ADD{havia recebido}} dois ganhou também outros dois.
\VS{18}Mas o que tinha recebido um foi cavar a terra, e escondeu o dinheiro do seu senhor.
\VS{19}Muito tempo depois, o senhor daqueles servos veio fazer contas com eles.
\VS{20}O que havia recebido cinco talentos chegou trazendo-lhe outros cinco talentos, e disse: ”Senhor, cinco talentos me entregaste, eis que ganhei com eles outros cinco talentos”.
\VS{21}E o seu senhor lhe disse: “
{\ADD{Muito}} bem, servo bom e fiel! Sobre o pouco foste fiel, sobre o muito te porei; entra na alegria do teu senhor”.
\VS{22}E chegando-se também o que havia recebido dois talentos, disse: “Senhor, dois talentos me entregaste, eis que ganhei com eles outros dois talentos”.
\VS{23}Seu senhor lhe disse: “
{\ADD{Muito}} bem, servo bom e fiel! Sobre o pouco foste fiel, sobre o muito te porei; entra na alegria do teu senhor”.
\VS{24}Mas, chegando também o que havia recebido um talento, disse: “Senhor, eu te conhecia, que és homem duro, que colhes onde não semeaste, e ajuntas onde não espalhaste;
\VS{25}E eu, atemorizado, fui e escondi o teu talento na terra; eis aqui tens o que é teu”.
\VS{26}Porém seu senhor lhe respondeu: “Servo mau e preguiçoso! Sabias que colho onde não semeei, e ajunto onde não espalhei.
\VS{27}Devias, portanto, ter depositado o meu dinheiro com os banqueiros e, quando eu voltasse, receberia o que é meu com juros.
\VS{28}Por isso, tirai dele o talento, e dai-o ao que tem dez talentos”.
\VS{29}Pois a todo aquele que tiver, lhe será dado, e terá em abundância; porém ao que não tiver, até o que tem lhe será tirado.
\VS{30}“E lançai o servo inútil às trevas de fora (ali haverá pranto e ranger de dentes)”.
\VS{31}E quando o Filho do homem vier em sua glória, e todos os santos anjos com ele, então ele se assentará sobre o trono de sua glória.
\VS{32}E serão ajuntadas diante dele todas as nações, e separará
{\ADD{as pessoas}} umas das outras, assim como o pastor separa as ovelhas dos bodes.
\VS{33}E porá as ovelhas à sua direita, porém os bodes à esquerda.
\VS{34}Então o Rei dirá o rei aos que estiverem à sua direita: “Vinde, benditos de meu Pai! Herdai o Reino que está preparado para vós desde a fundação do mundo.
\VS{35}Pois tive fome, e me destes de comer; tive sede, e me destes de beber; fui forasteiro, e me acolhestes;
\VS{36}{\ADD{estive}} nu, e me vestistes;
{\ADD{estive}} doente, e cuidastes de mim; estive na prisão, e me visitastes”.
\VS{37}Então os justos lhe perguntarão: “Senhor, quando te vimos com fome, e
{\ADD{te}} demos de comer, ou com sede, e
{\ADD{te}} demos de beber?
\VS{38}E quando te vimos forasteiro, e
{\ADD{te}} acolhemos, ou nu, e te vestimos?
\VS{39}E quando te vimos doente, ou na prisão, e viemos te visitar?”
\VS{40}E o Rei lhes responderá: “Em verdade vos digo que, todas as vezes que fizestes a um destes menores dos meus irmãos, fizestes a mim”.
\VS{41}Então dirá também aos que
{\ADD{estiverem}} à esquerda: “Apartai-vos de mim, malditos, ao fogo eterno, preparado para o Diabo e os seus anjos.
\VS{42}Pois tive fome, e não me destes de comer; tive sede, e não me destes de beber.
\VS{43}Fui forasteiro, e não me acolhestes;
{\ADD{estive}} nu, e não me vestistes;
{\ADD{estive}} doente, e na prisão, e não me visitastes.
\VS{44}Então também eles lhe perguntarão: “Senhor, quando te vimos com fome, ou com sede, ou forasteiro, ou nu, ou doente, ou na prisão, e não te servimos?”
\VS{45}Então ele lhes responderá, dizendo: “Em verdade vos digo que, todas as vezes que não fizestes a um destes menores, não fizestes a mim”.
\VS{46}E estes irão ao tormento eterno, porém os justos à vida eterna.

\par }\Chap{26}{\PP \VerseOne{1}E aconteceu que, quando Jesus terminou todas estas palavras, disse aos seus discípulos:
\VS{2}Vós bem sabeis que daqui a dois dias é a Páscoa, e o Filho do homem será entregue para ser crucificado.
\VS{3}Então os chefes dos sacerdotes, os escribas, e os anciãos do povo se reuniram na casa do sumo sacerdote, que se chamava Caifás.
\VS{4}E conversaram a fim de, usando mentira,
\FTNT{*}{{\FR 26:4 }
usando mentira
ou: em segredo
} prenderem Jesus, e
{\ADD{o}} matarem.
\VS{5}Porém diziam: Não na festa, para que não haja tumulto entre o povo.
\VS{6}Enquanto Jesus estava em Betânia, na casa de Simão o leproso,
\VS{7}veio a ele uma mulher com um vaso de alabastro, de óleo perfumado de grande valor, e derramou sobre a cabeça dele, enquanto estava sentado
\FTNT{†}{{\FR 26:7 }
sentado
ou: reclinado
} à mesa.
\VS{8}E quando os seus discípulos viram, ficaram indignados, dizendo: Para que este desperdício?
\VS{9}Pois esse óleo perfumado podia ter sido vendido por muito, e o dinheiro dado aos pobres.
\VS{10}Porém Jesus, sabendo
{\ADD{disso}} ,disse-lhes: Por que perturbais a esta mulher? Ora, ela me fez uma boa obra!
\VS{11}Pois vós sempre tendes os pobres convosco, porém nem sempre me tereis.
\VS{12}Pois ela, ao derramar este óleo perfumado sobre o meu corpo, ela o fez para
{\ADD{preparar}} o meu sepultamento.
\VS{13}Em verdade vos digo que, onde quer que este Evangelho em todo o mundo for pregado, também se dirá o que ela fez, para que seja lembrada.
\VS{14}Então um dos doze, chamado Judas Iscariotes, foi aos chefes dos sacerdotes,
\VS{15}e disse: O que quereis me dar, para que eu o entregue a vós? E eles lhe determinaram trinta
{\ADD{moedas}} de prata.
\VS{16}E desde então ele buscava oportunidade para o entregar.
\FTNT{‡}{{\FR 26:16 }
entregar
ou: trair
}
\VS{17}E no primeiro
{\ADD{dia da festa}} dos pães sem fermento, os discípulos vieram a Jesus lhe perguntar: Onde queres que te preparemos para comer a Páscoa?
\VS{18}E ele respondeu: Ide à cidade a um tal, e dizei-lhe: “O Mestre diz: ‘Meu tempo está perto. Contigo
\FTNT{§}{{\FR 26:18 }
Contigo
i.e. em tua casa
} celebrarei a Páscoa com os meus discípulos’ ”.
\VS{19}Os discípulos fizeram como Jesus havia lhes mandado, e prepararam a Páscoa.
\VS{20}E vindo o anoitecer, ele se assentou
\FTNT{*}{{\FR 26:20 }
assentou
ou: reclinou
} à mesa com os doze.
\VS{21}E enquanto comiam, disse: Em verdade vos digo que um de vós me trairá.
\FTNT{†}{{\FR 26:21 }
trairá
ou: entregará – também no v. 23
}
\VS{22}Eles ficaram muito tristes, e cada um deles começou a lhe perguntar: Por acaso sou eu, Senhor?
\VS{23}E ele respondeu: O que mete comigo a mão no prato, esse me trairá.
\VS{24}De fato, o Filho do homem vai assim como dele está escrito; mas ai daquele homem por quem o Filho do homem é traído! Bom seria a tal homem se não houvesse nascido.
\VS{25}E Judas, o que o traía, perguntou: Por acaso sou eu, Rabi?
{\ADD{Jesus}} lhe disse: Tu o disseste.
\VS{26}E enquanto comiam, Jesus tomou o pão, abençoou-o, e o partiu. Então o deu aos discípulos, e disse: Tomai, comei; isto é o meu corpo.
\VS{27}Em seguida tomou o cálice, deu graças, e o deu a eles, dizendo: Bebei dele todos,
\VS{28}porque este é o meu sangue, o
{\ADD{sangue}} do novo testamento,
\FTNT{‡}{{\FR 26:28 }
testamento
ou: pacto
} o qual é derramado por muitos, para o perdão de pecados.
\VS{29}E eu vos digo que desde agora não beberei deste fruto da vide, até aquele dia, quando convosco o beber, novo, no reino do meu Pai.
\VS{30}E depois de cantarem um hino,
\FTNT{§}{{\FR 26:30 }
hino
ou: salmo
} saíram para o monte das Oliveiras.
\VS{31}Então Jesus lhes disse: Todos vós falhareis comigo
\FTNT{*}{{\FR 26:31 }
falhareis comigo
Ou: me abandonareis. Tradicionalmente: “vós vos encandalizareis de mim”
} esta noite; porque está escrito: “Ferirei o pastor, e as ovelhas do rebanho serão dispersas”.
\VS{32}Mas, depois que eu for ressuscitado, irei adiante de vós para a Galileia.
\VS{33}Pedro, porém, respondeu-lhe: Ainda que todos falhem contigo, eu nunca falharei.
\VS{34}Jesus lhe disse: Em verdade te digo que, nesta mesma noite, antes do galo cantar, tu me negarás três vezes.
\VS{35}Pedro lhe respondeu: Ainda que eu tenha de morrer contigo, em nenhuma maneira te negarei. E todos os discípulos disseram o mesmo.
\VS{36}Então Jesus veio com eles a um lugar chamado Getsêmani, e disse aos discípulos: Ficai sentados aqui, enquanto eu vou ali orar.
\VS{37}Enquanto trazia consigo Pedro e os dois filhos de Zebedeu, ele começou a se entristecer e a se angustiar muito.
\VS{38}Então lhes disse: Minha alma está completamente triste até a morte. Ficai aqui, e vigiai comigo.
\VS{39}E indo um pouco mais adiante, prostrou-se sobre o seu rosto, orando, e dizendo: Meu Pai, se é possível, passe de mim este cálice; porém, não
{\ADD{seja}} como eu quero, mas sim como tu
{\ADD{queres}} .
\VS{40}Então voltou aos seus discípulos, e os encontrou dormindo; e disse a Pedro: Então, nem sequer uma hora pudestes vigiar comigo?
\VS{41}Vigiai e orai, para que não entreis em tentação. De fato, o espírito
{\ADD{está}} pronto, mas a carne
{\ADD{é}} fraca.
\VS{42}Ele foi orar pela segunda vez, dizendo: Meu Pai, se este cálice não pode passar de mim sem que eu o beba, faça-se a tua vontade.
\VS{43}Quando voltou, achou-os outra vez dormindo, pois os seus olhos estavam pesados.
\VS{44}Então os deixou, e foi orar pela terceira vez, dizendo as mesmas palavras.
\VS{45}Depois veio aos seus discípulos, e disse-lhes: Agora dormi e descansai.
\FTNT{†}{{\FR 26:45 }
Agora dormi, e descansai.
Ou: “Ainda dormis, e descansais?”
} Eis que chegou a hora em que o Filho do homem é entregue em mãos de pecadores.
\VS{46}Levantai-vos, vamos! Eis que chegou o que me trai.
\VS{47}Enquanto ele ainda estava falando, eis que veio Judas, um dos doze, e com ele uma grande multidão, com espadas e bastões, da parte dos chefes dos sacerdotes e dos anciãos do povo.
\VS{48}O seu traidor havia lhes dado sinal, dizendo: Aquele a quem eu beijar, é esse. Prendei-o.
\VS{49}Logo ele se aproximou de Jesus, e disse: Felicitações, Rabi! e o beijou.
\VS{50}Jesus, porém, lhe perguntou: Amigo, para que vieste? Então chegaram, agarraram Jesus,
\FTNT{‡}{{\FR 26:50 }
agarraram Jesus
lit. puseram as mãos em Jesus
} e o prenderam.
\VS{51}E eis que um dos que
{\ADD{estavam}} com Jesus estendeu a mão, puxou de sua espada, e feriu o servo do sumo sacerdote, cortando-lhe uma orelha.
\VS{52}Jesus, então, lhe disse: Põe de volta tua espada ao seu lugar, pois todos os que pegarem espada, pela espada perecerão.
\VS{53}Ou, por acaso, pensas tu que eu não posso agora orar ao meu Pai, e ele me daria mais de doze legiões de anjos?
\VS{54}Como, pois, se cumpririam as Escrituras
{\ADD{que dizem}} que assim tem que ser feito?
\VS{55}Naquela hora Jesus disse às multidões: Como a um ladrão saístes com espadas e bastões para me prender? Todo dia eu me sentava convosco, ensinando no templo, e não me prendestes.
\VS{56}Porém tudo isto aconteceu para que as Escrituras dos profetas se cumpram.Então todos os discípulos o abandonaram, e fugiram.
\VS{57}Os que prenderam Jesus o trouxeram
{\ADD{à casa}} de Caifás, o sumo sacerdote, onde os escribas e os anciãos estavam reunidos.
\VS{58}E Pedro o seguia de longe, até o pátio do sumo sacerdote; e entrou, e se assentou com os servos, para ver o fim.
\VS{59}Os chefes dos sacerdotes, os anciãos, e todo o supremo conselho
\FTNT{§}{{\FR 26:59 }
supremo conselho
lit. sinédrio – o conselho ou tribunal para assuntos da religião judaica
} buscavam falso testemunho contra Jesus, para poderem matá-lo,
\VS{60}mas não encontravam. E ainda que muitas falsas testemunhas se apresentavam,
{\ADD{contudo}} não encontravam.
\VS{61}Mas, por fim, vieram duas falsas testemunhas, que disseram: Este disse: “Posso derrubar o Templo de Deus e reconstruí-lo em três dias”.
\VS{62}Então o sumo sacerdote se levantou, e lhe perguntou: Não respondes nada ao que eles testemunham contra ti?
\VS{63}Porém Jesus ficava calado. Então o sumo sacerdote lhe disse: Ordeno-te
\FTNT{*}{{\FR 26:63 }
Ordeno-te
lit. conjuro-te, isto é, obrigar alguém a dizer algo sob juramento
} pelo Deus vivo que nos digas se tu és o Cristo, o Filho de Deus.
\VS{64}Jesus lhe disse: Tu o disseste. Porém eu vos digo que, desde agora, vereis o Filho do homem, sentado à direita do Poderoso,
\FTNT{†}{{\FR 26:64 }
Poderoso
lit. Poder
} e vindo sobre as nuvens do céu.
\VS{65}Então o sumo sacerdote rasgou suas roupas, e disse: Ele blasfemou! Para que necessitamos mais de testemunhas? Eis que agora ouvistes a sua blasfêmia.
\VS{66}Que vos parece?E eles responderam: Culpado de morte ele é.
\VS{67}Então lhe cuspiram no rosto, e lhe deram socos.
\VS{68}Outros lhe deram bofetadas, e diziam: Profetiza-nos, ó Cristo, quem é o que te feriu?
\VS{69}Pedro estava sentado fora no pátio. Uma serva aproximou-se dele, e disse: Também tu estavas com Jesus, o galileu.
\VS{70}Mas ele o negou diante de todos, dizendo: Não sei o que dizes.
\VS{71}E quando ele saiu em direção à entrada, outra o viu, e disse aos que ali
{\ADD{estavam}} : Também este estava com Jesus, o nazareno.
\VS{72}E ele o negou outra vez com um juramento: Não conheço
{\ADD{esse}} homem.
\VS{73}Pouco depois, os que ali estavam se aproximaram, e disseram a Pedro: Verdadeiramente também tu és um deles, pois a tua fala
\FTNT{‡}{{\FR 26:73 }
a tua fala
ou: “o teu sotaque”
} te denuncia.
\VS{74}Então ele começou a amaldiçoar e a jurar: Não conheço
{\ADD{esse}} homem!E imediatamente o galo cantou.
\VS{75}Então Pedro se lembrou da palavra de Jesus, que lhe dissera: Antes do galo cantar, tu me negarás três vezes. Assim ele saiu, e chorou amargamente.

\par }\Chap{27}{\PP \VerseOne{1}Vinda a manhã, todos os chefes dos sacerdotes e anciãos do povo juntamente se aconselharam contra Jesus, para o matarem.
\VS{2}E o levaram amarrado, e o entregaram a Pôncio Pilatos, o governador.
\VS{3}Então Judas, o que o havia traído, ao ver que
{\ADD{Jesus}} já estava condenado, devolveu, sentindo remorso, as trinta
{\ADD{moedas}} de prata aos chefes dos sacerdotes e aos anciãos;
\VS{4}e disse: Pequei, traindo sangue inocente. Porém eles disseram: Que nos interessa? Isso é problema teu!
\FTNT{*}{{\FR 27:4 }
Isso é problema teu
lit. Vê [isso] tu
}
\VS{5}Então ele lançou as
{\ADD{moedas}} de prata no templo, saiu, e foi enforcar-se.
\VS{6}Os chefes dos sacerdotes tomaram as
{\ADD{moedas}} de prata, e disseram: Não é lícito pô-las no tesouro das ofertas, pois isto é preço de sangue.
\VS{7}Então juntamente se aconselharam, e compraram com elas o campo do oleiro, para ser cemitério dos estrangeiros.
\VS{8}Por isso aquele campo tem sido chamado campo de sangue até hoje.
\VS{9}Assim se cumpriu o que foi dito pelo profeta Jeremias, que disse: Tomaram as trinta
{\ADD{moedas}} de prata, preço avaliado pelos filhos de Israel, o qual eles avaliaram;
\VS{10}e as deram pelo campo do oleiro, conforme o que o Senhor me mandou.
{\RQ{Zacarias 11:12-13; Jeremias 19:1-13; 32:6-9
}}
\VS{11}Jesus esteve diante do governador, e o governador lhe perguntou: És tu o Rei dos Judeus? E Jesus lhe respondeu: Tu
{\ADD{o}} dizes.
\VS{12}E, sendo ele foi acusado pelos chefes dos sacerdotes e pelos anciãos, nada respondeu.
\VS{13}Pilatos, então, lhe disse: Não ouves quantas coisas estão testemunhando contra ti?
\VS{14}Mas
{\ADD{Jesus}} não lhe respondeu uma só palavra, de maneira que o governador ficou muito maravilhado.
\VS{15}Na festa o governador costuma soltar um preso ao povo, qualquer um que quisessem.
\VS{16}E tinham então um preso bem conhecido, chamado Barrabás.
\VS{17}Quando, pois, se ajuntaram, Pilatos lhes perguntou: Qual quereis que vos solte? Barrabás, ou Jesus, que é chamado Cristo?
\VS{18}Pois ele sabia que foi por inveja que o entregaram.
\VS{19}E, enquanto ele estava sentado no assento de juiz, sua mulher lhe enviou a seguinte mensagem: Nada
{\ADD{faças}} com aquele justo, pois hoje sofri muito em sonhos por causa dele.
\VS{20}Mas os chefes dos sacerdotes e os anciãos persuadiram as multidões a pedirem Barrabás, e a exigirem a morte de Jesus.
\VS{21}O governador lhes perguntou: Qual destes dois quereis que vos solte? E responderam: Barrabás!
\VS{22}Pilatos lhes disse: Que, pois, farei de Jesus, que é chamado Cristo? Todos lhe disseram: Seja crucificado!
\VS{23}E o governador perguntou: Ora, que mal ele fez? Porém gritavam mais: Seja crucificado!
\VS{24}Quando, pois, Pilatos viu que nada adiantava, em vez disso se fazia mais tumulto, ele pegou água, lavou as mãos diante da multidão, e disse: Estou inocente do sangue deste justo. A responsabilidade é vossa.
\FTNT{†}{{\FR 27:24 }
a responsabilidade é vossa
lit. vós mesmos, vede
}
\VS{25}E todo o povo respondeu: O sangue dele
{\ADD{venha}} sobre nós, e sobre os nossos filhos.
\VS{26}Então soltou-lhes Barrabás, enquanto que mandou açoitar Jesus, e o entregou para ser crucificado.
\VS{27}Em seguida, os soldados do governador levaram Jesus consigo ao pretório, ajuntaram-se a ele toda a unidade miltar.
\FTNT{‡}{{\FR 27:27 }
unidade militar
ou: coorte, uma unidade de aproximadamente 500 soldados
}
\VS{28}Eles o despiram e o cobriram com um manto vermelho.
\VS{29}E, depois de tecerem uma coroa de espinhos, puseram-na sobre a sua cabeça, e uma cana em sua mão direita. Em seguida, puseram-se de joelhos diante dele, zombando-o, e diziam: Felicitações, Rei dos Judeus!
\VS{30}E cuspiram nele, tomaram a cana, e deram-lhe golpes na cabeça.
\VS{31}Depois de terem o zombado, despiram-lhe a capa, vestiram-no com suas roupas, e o levaram para crucificar.
\VS{32}Ao saírem, encontraram um homem de Cirene, por nome Simão; e obrigaram-no a levar sua cruz.
\VS{33}E quando chegaram ao lugar chamado Gólgota, que significa “o lugar da caveira”,
\VS{34}deram-lhe de beber vinagre misturado com fel. E, depois de provar, não quis beber.
\VS{35}E havendo-o crucificado, repartiram suas roupas, lançando sortes; para que se cumprisse o que foi dito pelo profeta: Repartiram entre si minhas roupas, e sobre minha túnica lançaram sortes.
\VS{36}Então se sentaram, e ali o vigiavam.
\VS{37}E puseram, por cima de sua cabeça, sua acusação escrita: ESTE É JESUS, O REI DOS JUDEUS.
\VS{38}Então foram crucificados com ele dois criminosos,
\FTNT{§}{{\FR 27:38 }
criminosos
ou: ladrões
} um à direita, e outro à esquerda.
\VS{39}Os que passavam blasfemavam dele, balançando suas cabeças,
\VS{40}e dizendo: Tu, que derrubas o Templo, e em três dias o reconstróis, salva a ti mesmo! Se és Filho de Deus, desce da cruz.
\VS{41}E da mesma maneira também os chefes dos sacerdotes, com os escribas e os anciãos, escarnecendo
{\ADD{dele}} ,diziam:
\VS{42}Salvou outros, a si mesmo não pode salvar. Se é Rei de Israel, desça agora da cruz, e creremos nele.
\VS{43}Confiou em Deus, livre-o agora, se lhe quer bem; pois disse: “Sou Filho de Deus”.
\VS{44}E os ladrões que estavam crucificados com ele também lhe insultavam.
\VS{45}Desde a hora sexta
\FTNT{*}{{\FR 27:45 }
hora sexta
aproximadamente meio-dia
} houve trevas sobre toda a terra até a hora nona.
\FTNT{†}{{\FR 27:45 }
hora nona
aproximadamente 3 horas da tarde
}
\VS{46}E perto da hora nona, Jesus gritou em alta voz: Eli, Eli, lamá sabactâni?, Isto é: Deus meu, Deus meu, porque me desamparaste?
\VS{47}E alguns dos que ali estavam, quando ouviram, disseram: Ele está chamando Elias.
\VS{48}Logo um deles correu e tomou uma esponja. Então a encheu de vinagre, colocou-a em uma cana, e lhe dava de beber.
\VS{49}Porém os outros disseram: Deixa, vejamos se Elias vem livrá-lo.
\VS{50}Jesus gritou outra vez em alta voz, e entregou o espírito.
\FTNT{‡}{{\FR 27:50 }
espírito
ou: fôlego, i.e. parou de respirar
}
\VS{51}E eis que o véu do Templo se rasgou em dois, de cima até embaixo, a terra tremeu, e as pedras se fenderam.
\VS{52}Os sepulcros se abriram, e muitos corpos de santos que tinham morrido
\FTNT{§}{{\FR 27:52 }
morrido
lit. dormido
} foram ressuscitados.
\VS{53}E, depois de ressuscitarem, saíram dos sepulcros, vieram à santa cidade,
\FTNT{*}{{\FR 27:53 }
santa cidade
i.e. Jerusalém
} e apareceram a muitos.
\VS{54}E o centurião, e os que com ele vigiavam Jesus, ao verem o terremoto e as coisas que haviam sucedido, tiveram muito medo, e disseram: Verdadeiramente ele era Filho de Deus.
\VS{55}Muitas mulheres, que desde a Galileia haviam seguido Jesus, e o serviam, estavam ali, olhando de longe.
\VS{56}Entre elas estavam Maria Madalena, e Maria mãe de Tiago e de José, e a mãe dos filhos de Zebedeu.
\VS{57}E chegado o entardecer, veio um homem rico de Arimateia, por nome José, que também era discípulo de Jesus.
\VS{58}Ele chegou a Pilatos, e pediu o corpo de Jesus. Então Pilatos mandou que o corpo
{\ADD{lhe}} fosse entregue.
\VS{59}José tomou o corpo, e o envolveu em um lençol limpo, de linho fino,
\VS{60}e o pôs em seu sepulcro novo, que tinha escavado numa rocha; em seguida rolou uma grande pedra à porta do sepulcro, e foi embora.
\VS{61}E ali estavam Maria Madalena e a outra Maria, sentadas de frente ao sepulcro.
\VS{62}No dia seguinte, que é o depois da preparação, os chefes dos sacerdotes, e os fariseus se reuniram com Pilatos,
\VS{63}e disseram: Senhor, nos lembramos que aquele enganador, enquanto ainda vivia, disse: “Depois de três dias serei ressuscitado”.
\VS{64}Portanto, manda que o sepulcro esteja em segurança até o terceiro dia, para que não aconteça dos seus discípulos virem de noite, e o furtem, e digam ao povo que ele ressuscitou dos mortos; e
{\ADD{assim}} o último engano será pior que o primeiro.
\VS{65}Pilatos lhes disse: Vós tendes uma guarda. Ide fazer segurança como o entendeis.
\VS{66}E eles se foram, e fizeram segurança no sepulcro com a guarda, selando a pedra.

\par }\Chap{28}{\PP \VerseOne{1}No fim do sábado, quando já começava a clarear para o primeiro dia da semana, Maria Madalena, e a outra Maria vieram ver o sepulcro.
\VS{2}E eis que houve um grande terremoto; porque um anjo do Senhor desceu do céu, chegou, e moveu a pedra da entrada, e ficou sentado sobre ela.
\VS{3}A aparência dele era como um relâmpago, e sua roupa branca como neve.
\VS{4}E de medo dele os guardas tremeram muito, e ficaram como mortos.
\VS{5}Mas o anjo disse às mulheres: Não vos atemorizeis, pois eu sei que buscais Jesus, o que foi crucificado.
\VS{6}Ele não está aqui, pois já ressuscitou, como ele disse. Vinde ver o lugar onde o Senhor jazia.
\VS{7}Ide depressa dizer aos seus discípulos que ele ressuscitou dos mortos; e eis que vai adiante de vós para a Galileia; ali o vereis. Eis que eu tenho vos dito.
\VS{8}Então elas saíram apressadamente do sepulcro, com temor e grande alegria, e correram para anunciar aos seus discípulos.
\VS{9}E, enquanto elas iam anunciar aos seus discípulos, eis que Jesus veio ao encontro delas, e disse: Saudações. Elas se aproximaram, pegaram os pés dele, e o adoraram.
\VS{10}Jesus, então, lhes disse: Não temais. Ide anunciar aos meus irmãos para eles irem à Galileia, e ali me verão.
\VS{11}Enquanto elas iam, eis que alguns da guarda vieram à cidade, e anunciaram aos chefes dos sacerdotes tudo o que havia acontecido.
\VS{12}Então eles se reuniram com os anciãos, depois de decidirem em conjunto, deram muito dinheiro aos soldados,
\VS{13}dizendo: Falai: “Os discípulos dele vieram de noite, e o furtaram enquanto estávamos dormindo”.
\VS{14}E, se isto for ouvido pelo governador, nós o persuadiremos, e vos manteremos seguros.
\VS{15}Eles tomaram o dinheiro e fizeram como foram instruídos. E este dito foi divulgado entre os judeus até hoje.
\VS{16}Os onze discípulos se foram para a Galileia, ao monte onde Jesus havia lhes ordenado.
\VS{17}E quando o viram, o adoraram; porém alguns duvidaram.
\VS{18}Jesus se aproximou deles, e lhes falou: Todo o poder me é dado no céu e na terra.
\VS{19}Portanto ide, fazei discípulos a todas as nações, batizando-os em nome do Pai, do Filho, e do Espírito Santo,
\VS{20}ensinando-lhes a guardar todas as coisas que eu vos tenho mandado. E eis que eu estou convosco todos os dias, até o fim dos tempos.
\FTNT{*}{{\FR 28:20 }
\_ fim dos tempos : Ou: fim da era
} Amém.\par }